% ROZDZIAL 10 — Organizacja pamieci i wydajnosc w C++20
% Fragment do wlaczenia po rozdziale 9.
% Wymagane pakiety: listings, array, amsmath, hyperref.
% Uzycie: % ROZDZIAL 10 — Organizacja pamieci i wydajnosc w C++20
% Fragment do wlaczenia po rozdziale 9.
% Wymagane pakiety: listings, array, amsmath, hyperref.
% Uzycie: % ROZDZIAL 10 — Organizacja pamieci i wydajnosc w C++20
% Fragment do wlaczenia po rozdziale 9.
% Wymagane pakiety: listings, array, amsmath, hyperref.
% Uzycie: % ROZDZIAL 10 — Organizacja pamieci i wydajnosc w C++20
% Fragment do wlaczenia po rozdziale 9.
% Wymagane pakiety: listings, array, amsmath, hyperref.
% Uzycie: \input{rozdzial_10_pamiec_wydajnosc_cpp.tex}

\chapter[Organizacja pamięci i wydajność]{Organizacja pamięci, wyrównanie i wydajność niskopoziomowa w C++}
\label{chap:cpp-memory-performance}

\section{Po co analizować organizację pamięci?}

W rozdziałach o wskaźnikach (rozdział~3), semantyce przenoszenia
(rozdział~6) i iteratorach (rozdział~8) interesowało nas głównie
\emph{poprawne zarządzanie obiektami}. Teraz przyjrzymy się temu, jak
obiekty są reprezentowane w pamięci oraz w jaki sposób wybory
programisty wpływają na liczbę alokacji i lokalność dostępu do danych.

Rozróżniajmy trzy poziomy opisu:
\begin{enumerate}
\item \textbf{Gwarancje języka C++:} rozmiar obiektu, wymagane wyrównanie,
      czas życia obiektu, legalność dostępu i zasady alokowania pamięci.
\item \textbf{Szczegóły implementacji:} rozmieszczenie pól, rozmiary
      wskaźników, zachowanie konkretnego alokatora czy układ binarny ABI.
\item \textbf{Właściwości sprzętu:} linie cache, hierarchia pamięci,
      prefetching, przepustowość i zjawisko \emph{false sharing}.
\end{enumerate}

Kod poprawny według standardu może być wolny; szybki kod działający
na jednym komputerze może być powolny na innym. Dlatego optymalizacja
wymaga najpierw \emph{pomiaru}, a następnie weryfikacji hipotezy.

\section{Obiekt, bajt, adres i czas życia}
\label{sec:cpp10-object}

\subsection{Pamięć nie jest tym samym co obiekt}

W modelu C++ \emph{obszar pamięci} może istnieć bez żywego obiektu typu
\texttt{T}. Obiekt charakteryzują między innymi typ, rozmiar, czas życia
i, dla pewnych klas, odpowiednio wywołane konstruktory i destruktory.
Samo uzyskanie adresu wystarczająco dużego bloku nie oznacza, że
wolno traktować ten blok jak dowolny obiekt C++.

\begin{lstlisting}
#include <cstddef>
#include <type_traits>

static_assert(sizeof(char) == 1);
static_assert(sizeof(std::byte) == 1);
static_assert(sizeof(int) >= 1);

struct Point { double x, y; };
static_assert(sizeof(Point) >= 2 * sizeof(double));
\end{lstlisting}

\texttt{sizeof(T)} podaje rozmiar obiektu w \emph{bajtach C++}.
Jeden bajt nie musi mieć ośmiu bitów; ich liczbę określa
\texttt{CHAR\_BIT} z nagłówka \texttt{<climits>}.
Rozmiar typu obiektowego jest dodatni. Każdy obiekt ma też
wymagane \emph{wyrównanie}, które omówimy dalej.

\paragraph{Czas przechowywania a czas życia.}
Należy odróżnić \emph{storage duration} od \emph{object lifetime}:
\begin{itemize}
\item obiekt lokalny ma zazwyczaj automatyczny czas przechowywania;
\item obiekt z \texttt{static} ma statyczny czas przechowywania;
\item obiekt \texttt{thread\_local} ma czas przechowywania wątku;
\item pamięć otrzymana przez \texttt{new} lub alokator jest pamięcią
      o dynamicznym czasie przechowywania.
\end{itemize}
Nie należy utożsamiać tych kategorii z gwarantowanymi fizycznymi
obszarami „stos” i „sterta”: te nazwy opisują zwyczajowe implementacje,
a nie pełny model abstrakcyjny języka.

\begin{lstlisting}
#include <memory>

int global = 17;                 // static storage duration
thread_local int per_thread = 0; // thread storage duration

void demo() {
    int local = 5;               // automatic storage duration
    auto p = std::make_unique<int>(10);
    // Obiekt *p zyje w pamieci zarzadzanej dynamicznie.
}
\end{lstlisting}

\subsection{Reprezentacja obiektu i bezpieczne oglądanie bajtów}

C++ pozwala badać reprezentację bajtową obiektu przez
\texttt{char}, \texttt{unsigned char} lub \texttt{std::byte}.
Nie wolno jednak wyciągać z tego wniosku, że dowolny wskaźnik
\texttt{T*} można przekonwertować na \texttt{U*} i odczytać dane jako
\texttt{U}: takie postępowanie może naruszyć reguły wyrównania,
czasu życia i dostępu przez glvalue odpowiedniego typu.

\begin{lstlisting}
#include <array>
#include <cstddef>
#include <cstring>
#include <cstdint>

std::uint32_t x = 0x12345678;
std::array<std::byte, sizeof(x)> bytes{};
std::memcpy(bytes.data(), &x, sizeof(x));

std::uint32_t y = 0;
std::memcpy(&y, bytes.data(), sizeof(y));
// y == x
\end{lstlisting}

Kolejność bajtów zależy od \emph{endianness}. Można ją rozpoznać przez
\texttt{std::endian} (C++20, nagłówek \texttt{<bit>}), jednak do
serializacji najlepiej jawnie ustalać porządek bajtów, format pól
i wersję protokołu. Nie należy wysyłać całej struktury przez
\texttt{write(fd, \&object, sizeof(object))} jako przenośnego formatu:
padding, reprezentacja liczb i ABI mogą się różnić.

\subsection{\texttt{std::bit\_cast} i ograniczenia rzutowań}

\texttt{std::bit\_cast<To>(from)} (C++20, \texttt{<bit>}) tworzy wartość
\texttt{To} przez skopiowanie reprezentacji bitowej, o ile oba typy
są trywialnie kopiowalne i mają równy rozmiar. Dodatkowo reprezentacja
musi oznaczać poprawną wartość docelowego typu.

\begin{lstlisting}
#include <bit>
#include <cstdint>

float f = 1.0f;
// Przyklad zaklada sizeof(float) == sizeof(uint32_t).
// Na typowej platformie wynik to kod bitowy IEEE 754.
if constexpr (sizeof(float) == sizeof(std::uint32_t)) {
    auto bits = std::bit_cast<std::uint32_t>(f);
    (void)bits;
}
\end{lstlisting}

Nie mylmy \texttt{bit\_cast} z \texttt{reinterpret\_cast}:
pierwszy wytwarza \emph{wartość}, drugi może tylko zmienić sposób
interpretacji wskaźnika; sam nie tworzy obiektu docelowego.

\section{Layout struktury: \texttt{sizeof}, \texttt{alignof}, padding}

\begin{center}\small
\begin{tabularx}{\textwidth}{@{}p{0.33\textwidth}X@{}}
\toprule
Narzędzie & Co pokazuje lub zmienia \\
\midrule
\code{sizeof(T)} & Podaje rozmiar obiektu typu \code{T} w bajtach, wraz z ewentualnym paddingiem. \\
\code{alignof(T)} & Podaje wymagane wyrównanie adresu obiektu typu \code{T}. \\
\code{offsetof(T, field)} & Podaje przesunięcie składowej w typie o standardowym układzie. \\
\code{alignas(n)} & Żąda wyrównania co najmniej do granicy \code{n} dla obiektu lub typu. \\
\code{std::byte} & Reprezentuje jednostkę surowej pamięci bez znaczenia znaku tekstowego. \\
\bottomrule
\end{tabularx}
\end{center}
\label{sec:cpp10-layout}

\subsection{Skąd bierze się padding?}

Każdy typ obiektowy ma wymaganie wyrównania. Kompilator może wstawić
\emph{padding}, czyli nieużywane bajty, aby kolejne pola miały
odpowiednie adresy, a elementy tablicy struktur były poprawnie
wyrównane. Rozmieszczenie pól zależy od typu klasy i reguł języka;
konkretna liczba bajtów paddingu zależy od implementacji.

\begin{lstlisting}
#include <cstddef>
#include <iostream>

struct Unordered {
    char flag;
    double value;
    int count;
};

struct Ordered {
    double value;
    int count;
    char flag;
};

std::cout << sizeof(Unordered) << '\n';
std::cout << sizeof(Ordered) << '\n';
std::cout << alignof(Ordered) << '\n';
\end{lstlisting}

Na częstym ABI, przy \texttt{sizeof(double)==8},
\texttt{alignof(double)==8} i \texttt{sizeof(int)==4},
pierwszy typ zajmuje zwykle 24 bajty, a drugi 16. \textbf{Nie są to
wartości wymagane przez standard}. Samo przestawienie pól może
zmniejszyć zapotrzebowanie na pamięć milionów rekordów, ale może też
pogorszyć czytelność, stabilność ABI lub lokalność dostępu.

\begin{center}
\small
\begin{tabular}{|l|p{4.4cm}|p{4.4cm}|}
\hline
 & \textbf{Unordered (przykładowo)} & \textbf{Ordered (przykładowo)} \\
\hline
Początek & \texttt{flag}: 1 B & \texttt{value}: 8 B \\
\hline
Dalej & padding: 7 B & \texttt{count}: 4 B \\
\hline
Dalej & \texttt{value}: 8 B & \texttt{flag}: 1 B \\
\hline
Koniec & \texttt{count}: 4 B + padding: 4 B & padding: 3 B \\
\hline
Razem & 24 B & 16 B \\
\hline
\end{tabular}
\end{center}

\subsection{\texttt{offsetof}, standard-layout i ABI}

Dla typów \emph{standard-layout} makro \texttt{offsetof} pozwala
poznać przesunięcie pola od początku obiektu. Używanie go z innymi
typami klas od C++17 jest warunkowo wspierane, dlatego do analiz
przenośnych przyjmujemy tutaj standard-layout.

\begin{lstlisting}
#include <cstddef>
#include <cstdint>
#include <type_traits>

struct Header {
    std::uint32_t id;
    std::uint16_t version;
};
static_assert(std::is_standard_layout_v<Header>);
static_assert(offsetof(Header, version) >=
              sizeof(std::uint32_t));
\end{lstlisting}

W analizie ABI trzeba pamiętać o polach wirtualnych, dziedziczeniu,
przestrzeniach adresowych i konwencjach wywołania. Zmiana kolejności
pól publicznej struktury używanej w bibliotekach binarnych może
zmienić jej ABI nawet wtedy, gdy program źródłowy nadal się kompiluje.

\subsection{\texttt{alignof} i \texttt{alignas}}

\texttt{alignof(T)} zwraca wyrównanie typu \texttt{T};
\texttt{alignas(N)} pozwala narzucić silniejsze wyrównanie.
Należy zachować ograniczenia języka dotyczące dopuszczalnych wartości
\texttt{N}. Wyrównanie może zwiększyć rozmiar obiektu, co jest
istotne przy tablicach.

\begin{lstlisting}
#include <atomic>
#include <cstdint>

struct alignas(64) AlignedCounter {
    std::atomic<std::uint64_t> value{0};
};
static_assert(alignof(AlignedCounter) >= 64);
\end{lstlisting}

Wyrównanie do 64 bajtów ma sens tylko wtedy, gdy wynika z wymagań
sprzętu lub pomiarów. Nie jest to uniwersalna wielkość linii cache.
Ponadto \texttt{alignas} nie czyni dostępu do obiektu atomowym,
nie zastępuje \texttt{std::atomic} i samo nie zapewnia synchronizacji.

\subsection{Czy warto „pakować” struktury?}

Atrybuty kompilatora w rodzaju \texttt{\_\_attribute\_\_((packed))}
lub \texttt{\#pragma pack} nie stanowią przenośnych mechanizmów
standardowego C++20. Zmniejszenie rozmiaru może wymusić niewyrównany
dostęp, obniżyć wydajność, a na niektórych architekturach być
problemem poprawności. W kodzie protokołów lepiej często dekodować
pola jawnie z bufora bajtów i nie wiązać formatu z layoutem klasy.

\section{Wyrównane alokacje i czas życia obiektów}
\label{sec:cpp10-lifetime}

\subsection{Różnica między alokacją a konstruowaniem}

Klasyczne \texttt{new T(args)} wykonuje logicznie dwa zadania:
pozyskuje odpowiednio wyrównaną pamięć i inicjalizuje w niej obiekt.
Analogicznie \texttt{delete} niszczy obiekt i zwalnia pamięć.
Dla narzędzi niskopoziomowych warto rozdzielić te operacje.

\begin{lstlisting}
#include <new>

struct Widget { explicit Widget(int) {} };

void* mem = ::operator new(sizeof(Widget));
Widget* w = nullptr;
try {
    w = ::new (mem) Widget(5); // placement new
} catch (...) {
    ::operator delete(mem);
    throw;
}
w->~Widget();                 // koniec zycia obiektu
::operator delete(mem);      // zwolnienie samej pamieci
\end{lstlisting}

Taka konstrukcja jest poprawna dla \texttt{Widget} o wyrównaniu
obsługiwanym przez zwykłe \texttt{operator new}. Dla typów
\emph{over-aligned} należy korzystać z odpowiedniej alokacji
obsługującej \texttt{std::align\_val\_t} lub ze standardowego
alokatora, który uwzględnia wymagania wyrównania.

\paragraph{Dlaczego to nie jest zalecany idiom aplikacyjny?}
Obsługa wyjątków, wielokrotnych destruktorów, rozmiaru tablicy
i zgodności alokacji z dealokacją łatwo prowadzi do błędów.
W większości kodu wystarczają \texttt{std::vector},
\texttt{std::make\_unique} i obiekty automatyczne. Rozdzielanie
alokacji od konstruowania ma sens głównie przy implementowaniu
kontenerów, pul pamięci i innych abstrakcji zasobów.

\subsection{\texttt{std::construct\_at} i \texttt{std::destroy\_at} (C++20)}

Od C++20 można użyć funkcji z nagłówka \texttt{<memory>} do
jawnego rozpoczęcia i zakończenia czasu życia obiektu w odpowiednio
przygotowanej pamięci.

\begin{lstlisting}
#include <cstddef>
#include <memory>

struct Tracked {
    int value;
    explicit Tracked(int n) : value(n) {}
};

alignas(Tracked) std::byte buffer[sizeof(Tracked)];
Tracked* obj = std::construct_at(
    reinterpret_cast<Tracked*>(buffer), 42);
// ... korzystamy z *obj ...
std::destroy_at(obj);
\end{lstlisting}

\textbf{Ważne:} \texttt{std::destroy\_at} kończy czas życia obiektu, ale
nie zwalnia jego pamięci. Nie wolno wykonywać dwóch destrukcji tego
samego żywego obiektu ani odwoływać się do obiektu po destrukcji.
Ponadto rozmiar i wyrównanie bufora muszą odpowiadać konstruowanemu
typowi; sam \texttt{reinterpret\_cast} nie naprawia niewłaściwego
wyrównania.

\subsection{Wyrównany \texttt{operator new}}

Dla typów o szczególnie silnym wyrównaniu istnieją przeciążenia
alokacji i dealokacji wykorzystujące \texttt{std::align\_val\_t}
(C++17). Zwykłe \texttt{new Overaligned{}} wybiera odpowiedni
mechanizm automatycznie.

\begin{lstlisting}
#include <new>

struct alignas(128) Block { char data[128]; };
Block* b = new Block{};
delete b;
\end{lstlisting}

Jeśli ręcznie używamy \texttt{::operator new} z argumentem
\texttt{std::align\_val\_t\{128\}}, musimy dobrać zgodną operację
\texttt{::operator delete} z parametrem wyrównania. Mieszanie tych
rodzin funkcji jest błędem.

\subsection{Ponowne użycie tej samej pamięci i \texttt{std::launder}}

W tej samej pamięci można zakończyć życie jednego obiektu, a potem
rozpocząć życie kolejnego, ale dostęp przez wcześniej utworzone
wskaźniki podlega szczegółowym regułom języka. \texttt{std::launder}
(C++17) ma zastosowanie w określonych sytuacjach, gdy nowy obiekt
nie jest przezroczystym zamiennikiem poprzedniego; nie jest
uniwersalnym remedium na aliasing, nieprawidłowy czas życia
czy źle wyrównaną pamięć.

\paragraph{Reguła praktyczna.}
W kodzie aplikacyjnym nie sięgaj po \texttt{std::launder} bez
konkretnego problemu z czasem życia i bez zrozumienia odpowiednich
reguł. Preferuj własność zasobu wyrażoną typem RAII.

\section{Alokatory: pamięć, obiekty i bezpieczeństwo wyjątkowe}
\label{sec:cpp10-allocators}

\subsection{Standardowy model alokatora}

Kontenery standardowe używają \emph{alokatorów} do uzyskiwania
pamięci na swoje elementy. Alokator nie jest tym samym co funkcja
\texttt{new} i nie musi budować obiektu w chwili rezerwacji pamięci.
\texttt{std::allocator\_traits} zapewnia jednolity interfejs do
alokowania, konstruowania i niszczenia obiektów.

\begin{lstlisting}
#include <memory>

struct Item { explicit Item(int v) : value(v) {} int value; };
std::allocator<Item> alloc;
Item* p = alloc.allocate(2); // pamiec na dwa obiekty

std::allocator_traits<decltype(alloc)>::construct(alloc, p, 5);
std::allocator_traits<decltype(alloc)>::construct(alloc, p + 1, 7);

// teraz dwa obiekty sa zywe
std::allocator_traits<decltype(alloc)>::destroy(alloc, p + 1);
std::allocator_traits<decltype(alloc)>::destroy(alloc, p);
alloc.deallocate(p, 2);
\end{lstlisting}

Przykład dla przejrzystości zakłada, że oba konstruktory nie rzucają
wyjątków. \textbf{Pełna implementacja kontenera} musi zliczać
pomyślnie skonstruowane obiekty i w razie wyjątku zniszczyć tylko je,
a potem zwolnić pamięć. Właśnie z tego powodu lepiej opierać się na
gotowych kontenerach niż ręcznie pisać alokację elementów.

\subsection{Rezerwacja pojemności i realokacja \texttt{vector}}

\texttt{std::vector} przechowuje elementy w pamięci ciągłej, dzięki
czemu zapewnia dobrą lokalność dostępu. \texttt{reserve} może
ograniczyć liczbę realokacji, jeśli znamy przewidywaną wielkość
kontenera.

\begin{lstlisting}
#include <vector>

std::vector<int> values;
values.reserve(10'000);
for (int i = 0; i < 10'000; ++i)
    values.push_back(i);
\end{lstlisting}

\texttt{reserve(n)} nie zmienia \texttt{size()}, lecz zapewnia
\texttt{capacity() >= n}, chyba że zgłoszony zostanie wyjątek.
Realokacja unieważnia wskaźniki, referencje i iteratory do elementów.
Nie jest dobrym pomysłem wywoływanie \texttt{reserve(size()+1)}
przy każdym dodaniu elementu: można w ten sposób utracić
korzystną politykę geometrycznego wzrostu pojemności.

\subsection{Dlaczego własny alokator to nie pierwszy krok?}

Własny alokator w szablonie kontenera wiąże się z wymaganiami
dotyczącymi propagowania przy kopiowaniu, przenoszeniu i zamianie
kontenerów oraz z kwestią porównywania alokatorów.
Dla większości zastosowań wygodniejsza jest warstwa
\texttt{std::pmr}, którą poznamy teraz.

\section{\texttt{std::pmr}: polimorficzne zasoby pamięci}

\begin{center}\small
\begin{tabularx}{\textwidth}{@{}p{0.37\textwidth}X@{}}
\toprule
Mechanizm & Zastosowanie i ograniczenie \\
\midrule
\code{pmr::memory_resource} & Interfejs strategii przydziału i zwalniania pamięci używany przez alokatory PMR. \\
\code{pmr::vector<T>} & Kontener, którego alokator korzysta z wybranego zasobu pamięci. \\
\code{monotonic_buffer_resource} & Przydziela kolejne bloki szybko; pojedyncze zwolnienia nie odzyskują miejsca, zasób zwalnia je zbiorczo. \\
\code{unsynchronized_pool_resource} & Pula małych bloków bez synchronizacji między wątkami. \\
\code{synchronized_pool_resource} & Pula z synchronizacją, odpowiednia do współdzielenia między wątkami. \\
\bottomrule
\end{tabularx}
\end{center}
\label{sec:cpp10-pmr}

\subsection{Alokator jako decyzja w czasie działania}

\texttt{std::pmr} (C++17, \texttt{<memory\_resource>}) oddziela
\emph{rodzaj kontenera} od \emph{sposobu przydzielania pamięci}.
\texttt{std::pmr::vector<T>} jest aliasem \texttt{std::vector<T, ...>}
z alokatorem polimorficznym. Zasób \texttt{memory\_resource}
można wybierać w czasie działania programu.

\begin{lstlisting}
#include <memory_resource>
#include <vector>

std::pmr::monotonic_buffer_resource arena;
std::pmr::vector<int> numbers{&arena};
for (int i = 0; i < 100; ++i)
    numbers.push_back(i);
\end{lstlisting}

Kontener przechowuje \emph{wskaźnik do zasobu}, ale nie przejmuje
własności zasobu pamięci. Zasób musi
pozostać żywy wystarczająco długo, aby obsłużyć wszystkie powiązane
alokacje i dealokacje. \textbf{Nie wolno niszczyć lub resetować
areny, gdy żywe obiekty wciąż korzystają z jej pamięci}.

\subsection{Arena monotoniczna: wiele alokacji, jedno zwolnienie}

\texttt{std::pmr::monotonic\_buffer\_resource} wydaje kolejne
fragmenty pamięci. Jej \texttt{deallocate} nie zwalnia pojedynczych
bloków; odzyskanie następuje przy \texttt{release()} lub destrukcji
zasobu. To pasuje do wielu obiektów budowanych razem i usuwanych
w jednej fazie, np. drzewa parsowania pojedynczego dokumentu.

\begin{lstlisting}
#include <cstddef>
#include <memory_resource>
#include <vector>

alignas(std::max_align_t) std::byte storage[4096];
std::pmr::monotonic_buffer_resource arena(
    storage, sizeof(storage),
    std::pmr::null_memory_resource());

{
    std::pmr::vector<int> data{&arena};
    for (int i = 0; i < 32; ++i) data.push_back(i);
} // kontener znika; jego deallocate nie uwalnia blokow
// arena jest teraz bezpieczna do zniszczenia lub release()
\end{lstlisting}

\texttt{null\_memory\_resource()} wymusza niepowodzenie alokacji,
gdy lokalny bufor się wyczerpie (rzucane jest \texttt{std::bad\_alloc}).
Jeżeli zamiast niego użyjemy zwykłego zasobu nadrzędnego
(\emph{upstream}), arena może dobierać kolejne bufory.

\paragraph{Ważna konsekwencja.}
Gdy \texttt{vector} zwiększa pojemność, porzuca stary blok.
W przypadku areny monotonicznej ten blok zwykle pozostaje zajęty
aż do końca całej fazy. Arena nie oznacza zatem automatycznie
najmniejszego zużycia pamięci; warto stosować \texttt{reserve}.

\subsection{Zasoby pulowe: \texttt{unsynchronized\_pool\_resource}}

Gdy obiekty są często alokowane i usuwane w różnych momentach,
przydatne bywają pule przechowujące bloki o podobnych rozmiarach.

\begin{lstlisting}
#include <memory_resource>
#include <vector>

std::pmr::unsynchronized_pool_resource pool;
std::pmr::vector<int> items{&pool};
items.resize(500);
\end{lstlisting}

\texttt{unsynchronized\_pool\_resource} nie zapewnia bezpieczeństwa
jednoczesnego użycia z wielu wątków. Do współdzielonego użycia
istnieje \texttt{synchronized\_pool\_resource}. Wybór zależy od
profilu alokacji i pomiarów: pule nie są uniwersalnie szybsze.

\subsection{Własny zasób: liczymy rzeczywiste alokacje}

Wariant diagnostyczny może zliczać żądania pamięci i delegować je
do \texttt{new\_delete\_resource()}. Trzy wirtualne funkcje
\texttt{do\_allocate}, \texttt{do\_deallocate}, \texttt{do\_is\_equal}
tworzą podstawę własnej implementacji.

\begin{lstlisting}
#include <cstddef>
#include <memory_resource>

struct CountingResource : std::pmr::memory_resource {
    std::pmr::memory_resource* upstream =
        std::pmr::new_delete_resource();
    std::size_t allocations = 0;
    std::size_t bytes = 0;
private:
    void* do_allocate(std::size_t n,
                      std::size_t alignment) override {
        void* p = upstream->allocate(n, alignment);
        ++allocations;
        bytes += n;
        return p;
    }
    void do_deallocate(void* p, std::size_t n,
                       std::size_t alignment) override {
        upstream->deallocate(p, n, alignment);
    }
    bool do_is_equal(
        const std::pmr::memory_resource& other
    ) const noexcept override {
        return this == &other;
    }
};
\end{lstlisting}

\begin{lstlisting}
CountingResource stats;
{
    std::pmr::vector<int> data{&stats};
    for (int i = 0; i < 100; ++i) data.push_back(i);
}
// stats.allocations i stats.bytes opisuja zadania
// przeslane do tego konkretnego zasobu.
\end{lstlisting}

\textbf{Uwaga:} nie są to wszystkie alokacje programu, a
\texttt{bytes} oznacza sumę żądanych rozmiarów, nie największe
jednoczesne zużycie pamięci. Warto liczyć także dealokacje,
\emph{peak usage} i rozkład rozmiarów bloków.

\subsection{Zagnieżdżone kontenery i propagacja zasobu}

\begin{lstlisting}
#include <memory_resource>
#include <string>
#include <vector>

CountingResource stats;
std::pmr::vector<std::pmr::string> names{&stats};
names.emplace_back("alpha");
names.emplace_back("a rather long text that needs allocation");
\end{lstlisting}

\texttt{pmr::string} również korzysta z polimorficznego alokatora,
ale dokładne zachowanie małych napisów zależy od implementacji
(\emph{small string optimization}). Nie należy zakładać, że każdy
\texttt{emplace\_back} przydziela pamięć ani że krótki napis wymaga
alokacji dynamicznej.

\section{Lokalność danych i hierarchia cache}
\label{sec:cpp10-cache}

\subsection{Dlaczego asymptotyka nie wystarcza?}

Dwa algorytmy o złożoności \(O(n)\) mogą mieć bardzo różną wydajność,
jeżeli pierwszy czyta sąsiadujące elementy tablicy, a drugi podąża
za wskaźnikami do rozproszonych obiektów. Procesor pobiera dane
porcjami (\emph{cache lines}); dostęp sekwencyjny sprzyja
prefetchingowi, zaś losowy dostęp bywa ograniczony opóźnieniem pamięci.

Nie ma jednego rozmiaru linii cache dla wszystkich procesorów ani
jednej wartości kosztu trafienia i chybienia. Zależą one od sprzętu,
aktualnego obciążenia i rodzaju dostępu.

\begin{lstlisting}
#include <numeric>
#include <vector>

std::vector<int> v(100000, 1);
long long sum = std::accumulate(v.begin(), v.end(), 0LL);
// Dostep sekwencyjny do pamieci ciaglej.
\end{lstlisting}

Wskaźniki do osobno alokowanych małych węzłów mogą pogorszyć
lokalność i powiększyć narzut pamięciowy. Nie oznacza to,
że \texttt{std::vector} jest zawsze lepszy od \texttt{std::list}:
wybór zależy od wymaganych operacji, unieważniania iteratorów
i rzeczywistego profilu obciążenia.

\subsection{AoS kontra SoA}

W układzie \emph{Array of Structures} (AoS) przechowujemy obok siebie
wszystkie składowe kolejnych rekordów. W \emph{Structure of Arrays}
(SoA) przechowujemy osobne tablice poszczególnych składowych.

\begin{lstlisting}
#include <vector>

struct PointAoS { float x, y, z, intensity; };
std::vector<PointAoS> aos;

struct PointSoA {
    std::vector<float> x, y, z, intensity;
};
PointSoA soa;
\end{lstlisting}

Gdy obliczamy wyłącznie sumę współrzędnych \texttt{x}, układ SoA
umożliwia odczytywanie samych potrzebnych danych, a nie całych
rekordów. Ułatwia to też wektoryzację. Jeśli natomiast w każdej
operacji potrzebujemy wszystkich pól jednego rekordu, AoS może być
korzystniejszy. Wydajność zależy od rozmiaru danych, kodu i sprzętu.

\begin{center}
\small
\begin{tabular}{|p{6.0cm}|p{6.0cm}|}
\hline
\textbf{AoS: rekordy obok siebie} & \textbf{SoA: kolumny obok siebie} \\
\hline
\texttt{[x1 y1 z1 i1] [x2 y2 z2 i2]} &
\texttt{[x1 x2 ...] [y1 y2 ...]} \\
\hline
Wygodny dostęp do całego punktu & Wygodny dostęp do jednej cechy wszystkich punktów \\
\hline
Często mniej kontenerów i prostszy model & Często lepsze wykorzystanie cache w obliczeniach kolumnowych \\
\hline
\end{tabular}
\end{center}

\subsection{Stride, skoki i przewidywanie dostępu}

Przechodzenie po pamięci ze skokiem \(k\) oznacza, że kolejne
operacje mogą dotykać odległych linii cache. Dla macierzy w układzie
wierszowym zwykle korzystniej jest przechodzić po kolumnach
\emph{wewnątrz pojedynczego wiersza}, a następnie po wierszach.
Zanim zmienimy kolejność pętli, należy uwzględnić zależności
obliczeniowe, autowektoryzację i to, czy kompilator już dokonał
danej transformacji.

\begin{lstlisting}
#include <vector>

constexpr std::size_t rows = 100;
constexpr std::size_t cols = 200;
std::vector<double> matrix(rows * cols, 1.0);
double sum = 0.0;
for (std::size_t i = 0; i < rows; ++i)
    for (std::size_t j = 0; j < cols; ++j)
        sum += matrix[i * cols + j];
\end{lstlisting}

\section{False sharing i współbieżne liczniki}
\label{sec:cpp10-false-sharing}

\subsection{Na czym polega false sharing?}

Jeżeli dwa wątki wielokrotnie zapisują \emph{różne} zmienne, ale
zmienne znajdują się w tej samej linii cache, protokół spójności
pamięci podręcznej może wymuszać kosztowną wymianę własności tej
linii między rdzeniami. Nazywamy to \emph{false sharing}.
Nie oznacza to \emph{data race}, jeżeli poszczególne zmienne są
atomowe lub inaczej poprawnie synchronizowane.

\begin{lstlisting}
#include <atomic>
#include <cstdint>

struct NaiveCounters {
    std::atomic<std::uint64_t> first{0};
    std::atomic<std::uint64_t> second{0};
};
\end{lstlisting}

Dwa liczniki mogą być blisko siebie, jednak dokładny layout obiektu
i mapowanie do linii cache nie są gwarantowane. Na niektórych
platformach ich równoczesna modyfikacja będzie kosztowna.

\subsection{Rozdzielenie liczników w pamięci}

\begin{lstlisting}
#include <atomic>
#include <cstdint>

struct alignas(64) IsolatedAtomic {
    std::atomic<std::uint64_t> value{0};
};
struct SeparatedCounters {
    IsolatedAtomic first, second;
};
\end{lstlisting}

Tu użyliśmy 64 bajtów wyłącznie jako \emph{przykładowej} wartości.
Od C++17 w nagłówku \texttt{<new>} istnieją dwie stałe
opisujące zalecane odstępy dla danych często modyfikowanych przez
różne wątki (\emph{destructive interference}) oraz dla danych,
które powinny pozostawać blisko siebie (\emph{constructive
interference}). Pełną nazwę pierwszej z nich pokazuje kod poniżej.
Wartości stałych zależą od implementacji; nie stanowią pomiaru
rozmiaru linii cache w uruchomionym procesorze.

\begin{lstlisting}
#include <atomic>
#include <new>

struct alignas(
    std::hardware_destructive_interference_size
) Counter {
    std::atomic<unsigned> value{0};
};
\end{lstlisting}

Rozdzielenie liczników zwiększa rozmiar struktury. Przy milionie
liczników koszt pamięciowy może przewyższyć korzyść. Technika ma
największy sens, gdy oddzielne wątki często zapisują niewielką
liczbę niezależnych pól.

\subsection{False sharing a semantyka atomowa}

\texttt{std::memory\_order\_relaxed} wystarczy do niektórych
niezależnych liczników, gdy zależy nam wyłącznie na atomowym
zliczaniu; nie gwarantuje jednak synchronizacji innych danych.
Usunięcie false sharing nie zastępuje poprawnego modelu pamięci
współbieżnej. Szczegółowe zasady operacji atomowych były omawiane
w rozdziale~1 o współbieżności.

\section{Jak rzetelnie mierzyć wydajność?}
\label{sec:cpp10-benchmark}

\subsection{Zasady pomiarów}

Najważniejsze pytanie brzmi nie „czy ten kod wygląda szybciej?”,
lecz „czy w rzeczywistym scenariuszu uzyskujemy powtarzalną poprawę?”.
Doświadczenie powinno uwzględniać:
\begin{enumerate}
\item kompilację zoptymalizowaną (np. \texttt{-O2} lub \texttt{-O3});
\item \emph{warm-up} i wielokrotne uruchomienia;
\item jednakowe dane oraz taki sam rezultat dla obu wariantów;
\item ochronę przed usunięciem testowanej pracy przez optymalizator;
\item pomiar rozkładu wyników, a nie pojedynczego odczytu;
\item ograniczenie wpływu alokacji, I/O, schedulerów i skalowania CPU;
\item kontrolę zużycia pamięci, a nie tylko czasu;
\item narzędzia profilujące (np. Linux \texttt{perf}, odpowiedniki
      sprzętowe lub profilery danego środowiska).
\end{enumerate}

\subsection{Minimalna funkcja pomiarowa}

\begin{lstlisting}
#include <chrono>
#include <utility>

template<class F>
auto measure(F&& f) {
    auto start = std::chrono::steady_clock::now();
    auto result = std::forward<F>(f)();
    auto stop = std::chrono::steady_clock::now();
    return std::pair{result, stop - start};
}
\end{lstlisting}

Taka funkcja przydaje się do nauki, ale \emph{nie jest kompletna}
biblioteka benchmarkowa. Dla zbyt krótkich fragmentów narzut
pomiaru może przeważać nad czasem pracy. Dla liczb zmiennoprzecinkowych
porównywanie różnych kolejności sumowania wymaga ostrożności,
ponieważ zmiana kolejności działań może zmienić zaokrąglenia.

\subsection{Sanitizatory i wydajność}

Sanitizatory, np. AddressSanitizer i UndefinedBehaviorSanitizer,
są narzędziami do badania poprawności pamięci i innych błędów.
Wprowadzają narzut, dlatego \emph{pomiarów porównawczych wydajności}
nie należy bezrefleksyjnie wykonywać na buildach sanitizowanych.
Bezpieczny cykl pracy wygląda następująco: test poprawności,
profilowanie zwykłego buildu, zmiana jednej rzeczy, ponowny test
poprawności i ponowny pomiar.

\section{Przykładowy projekt: szybkie przetwarzanie wielu punktów}
\label{sec:cpp10-project}

Przyjmijmy, że program przechowuje milion punktów \((x,y,z,I)\)
i dla każdego oblicza tylko prosty filtr zależny od \(x\) i \(I\).
Możemy rozważyć dwa modele:
\begin{enumerate}
\item \textbf{AoS:} \texttt{vector<Point>} --- prostszy interfejs,
      wygodny dostęp do pełnego punktu.
\item \textbf{SoA:} osobne \texttt{vector<float>} na \(x,y,z,I\) ---
      możliwy mniejszy transfer danych potrzebnych do filtru.
\end{enumerate}

Następnie oceniamy, czy warto użyć \texttt{reserve} przed zapełnieniem
kontenerów; jeśli rekordy są tworzone na potrzeby jednego zapytania,
możemy sprawdzić arenę monotoniczną. Trzeba jednak rozróżnić
\emph{koszt wypełnienia struktur} od \emph{kosztu przetwarzania}: arena
zmienia zachowanie alokacji, nie gwarantuje poprawy pracy procesora
nad danymi.

Dla kodu numerycznego duży wpływ może mieć też wybór kolejności
wymiarów macierzy, reprezentacji sąsiedztwa w siatce czy sposobu
obliczania operatora na konkretnym punkcie zamiast konstruowania
dużych obiektów pośrednich. Właściwa reprezentacja danych potrafi
dać większy zysk niż drobne optymalizacje pojedynczych instrukcji.

\section{Jak dobrać technikę?}

\begin{center}
\small
\begin{tabular}{|p{5.7cm}|p{6.7cm}|}
\hline
\textbf{Problem} & \textbf{Pierwszy sensowny krok} \\
\hline
Za dużo dynamicznych alokacji & \texttt{reserve}, ograniczenie obiektów pośrednich, pomiar alokacji \\
\hline
Wiele obiektów znika razem & \texttt{pmr::monotonic\_buffer\_resource} po kontroli czasu życia \\
\hline
Dużo małych alokacji o różnej żywotności & \texttt{pmr::...pool\_resource} i pomiar \\
\hline
Praca nad jedną kolumną wielu rekordów & Porównanie AoS z SoA \\
\hline
Losowe skoki po pamięci & Sprawdzenie lokalności, zmian formatu danych i blokowania pętli \\
\hline
Niezależne atomiki spowalniają się nawzajem & Analiza false sharing i separacja intensywnie zapisywanych pól \\
\hline
Trzeba uzyskać silne wyrównanie & \texttt{alignas}, odpowiednia alokacja i kontrola \texttt{alignof} \\
\hline
Trzeba napisać kontener lub pulę & \texttt{allocator\_traits}, zarządzanie czasem życia, RAII \\
\hline
Potrzeba przenośnej serializacji & Jawny format bajtowy, bez kopiowania całych struktur \\
\hline
\end{tabular}
\end{center}

\section{Typowe błędy i pułapki}

\begin{enumerate}
\item \textbf{Wniosek „\texttt{sizeof(T)} to suma pól”.}
      Padding może zwiększyć rozmiar, a reguły klas mogą być złożone.
\item \textbf{Uznanie layoutu z jednej kompilacji za standard.}
      ABI, wyrównanie i rozmiary typów różnią się między platformami.
\item \textbf{Nieprawidłowe \texttt{reinterpret\_cast}.}
      Sama zmiana typu wskaźnika nie legalizuje dostępu do danych.
\item \textbf{Mylenie pamięci z czasem życia obiektu.}
      \texttt{allocate} nie zastępuje konstruowania obiektów klasowych.
\item \textbf{Zwalnianie pamięci bez destruktora.}
      Dla nietrywialnych typów może wyciekać zasób zarządzany przez obiekt.
\item \textbf{Niezgodna para alokacja--dealokacja.}
      Szczególnie ważne przy over-alignment i ręcznej alokacji.
\item \textbf{Pozostawienie kontenera po zniszczeniu areny \texttt{pmr}.}
      Jego pamięć staje się niepoprawna; wynik to niezdefiniowane zachowanie.
\item \textbf{Wywołanie \texttt{release()} przy żywych obiektach.}
      Arena może unieważnić ich pamięć bez wywoływania destruktorów.
\item \textbf{Założenie, że \texttt{monotonic\_buffer\_resource} oszczędza
      pamięć.} Zwolnienie pojedynczych bloków zwykle nie następuje.
\item \textbf{Współdzielenie \texttt{unsynchronized\_pool\_resource}
      między wątkami bez ochrony.} Zasób nie jest thread-safe.
\item \textbf{Optymalizacja poprzez mniejszą strukturę bez pomiaru.}
      Mniejszy \texttt{sizeof} nie gwarantuje krótszego czasu wykonania.
\item \textbf{Rzekoma optymalizacja przez \texttt{volatile}.}
      \texttt{volatile} nie zapewnia synchronizacji między wątkami.
\item \textbf{Pomiar jednego przebiegu.} Scheduler, częstotliwości
      zegara i cache mogą silnie zniekształcić wynik.
\item \textbf{Użycie sanitizatora jako miary szybkości.}
      Narzut instrumentacji zmienia charakterystykę programu.
\end{enumerate}

\paragraph{Pytania kontrolne.}
\begin{enumerate}
\item Czym różnią się \texttt{sizeof(T)} i \texttt{alignof(T)}?
\item Dlaczego zmiana kolejności pól może zmienić rozmiar struktury?
\item Czy placement \texttt{new} przydziela pamięć?
\item Jaka jest różnica między \texttt{destroy\_at} a \texttt{deallocate}?
\item Kiedy arena monotoniczna może zużywać więcej pamięci niż oczekujemy?
\item Dlaczego \texttt{pmr::vector<T>} nie jest właścicielem zasobu?
\item Kiedy SoA może przyspieszyć obliczenia na rekordach?
\item Czy false sharing i data race są tym samym problemem?
\item Dlaczego \texttt{std::memory\_order\_relaxed} nie rozwiązuje
      wszystkich problemów synchronizacji?
\item Jak odróżnić realną poprawę wydajności od szumu pomiarowego?
\end{enumerate}

\paragraph{Dalsza lektura.}
Dokumentacja \emph{cppreference} obejmuje hasła \emph{Object},
\emph{Object lifetime}, \texttt{alignof}, \texttt{alignas},
\texttt{std::construct\_at} oraz \texttt{allocator\_traits}.
Osobno warto przeczytać opis zasobów \texttt{std::pmr}
i własności \emph{hardware interference size}.
W praktyce polecam analizę śladów alokacji, profilery
wydajności i \emph{Compiler Explorer} do porównywania
wygenerowanego kodu maszynowego.


\chapter[Organizacja pamięci i wydajność]{Organizacja pamięci, wyrównanie i wydajność niskopoziomowa w C++}
\label{chap:cpp-memory-performance}

\section{Po co analizować organizację pamięci?}

W rozdziałach o wskaźnikach (rozdział~3), semantyce przenoszenia
(rozdział~6) i iteratorach (rozdział~8) interesowało nas głównie
\emph{poprawne zarządzanie obiektami}. Teraz przyjrzymy się temu, jak
obiekty są reprezentowane w pamięci oraz w jaki sposób wybory
programisty wpływają na liczbę alokacji i lokalność dostępu do danych.

Rozróżniajmy trzy poziomy opisu:
\begin{enumerate}
\item \textbf{Gwarancje języka C++:} rozmiar obiektu, wymagane wyrównanie,
      czas życia obiektu, legalność dostępu i zasady alokowania pamięci.
\item \textbf{Szczegóły implementacji:} rozmieszczenie pól, rozmiary
      wskaźników, zachowanie konkretnego alokatora czy układ binarny ABI.
\item \textbf{Właściwości sprzętu:} linie cache, hierarchia pamięci,
      prefetching, przepustowość i zjawisko \emph{false sharing}.
\end{enumerate}

Kod poprawny według standardu może być wolny; szybki kod działający
na jednym komputerze może być powolny na innym. Dlatego optymalizacja
wymaga najpierw \emph{pomiaru}, a następnie weryfikacji hipotezy.

\section{Obiekt, bajt, adres i czas życia}
\label{sec:cpp10-object}

\subsection{Pamięć nie jest tym samym co obiekt}

W modelu C++ \emph{obszar pamięci} może istnieć bez żywego obiektu typu
\texttt{T}. Obiekt charakteryzują między innymi typ, rozmiar, czas życia
i, dla pewnych klas, odpowiednio wywołane konstruktory i destruktory.
Samo uzyskanie adresu wystarczająco dużego bloku nie oznacza, że
wolno traktować ten blok jak dowolny obiekt C++.

\begin{lstlisting}
#include <cstddef>
#include <type_traits>

static_assert(sizeof(char) == 1);
static_assert(sizeof(std::byte) == 1);
static_assert(sizeof(int) >= 1);

struct Point { double x, y; };
static_assert(sizeof(Point) >= 2 * sizeof(double));
\end{lstlisting}

\texttt{sizeof(T)} podaje rozmiar obiektu w \emph{bajtach C++}.
Jeden bajt nie musi mieć ośmiu bitów; ich liczbę określa
\texttt{CHAR\_BIT} z nagłówka \texttt{<climits>}.
Rozmiar typu obiektowego jest dodatni. Każdy obiekt ma też
wymagane \emph{wyrównanie}, które omówimy dalej.

\paragraph{Czas przechowywania a czas życia.}
Należy odróżnić \emph{storage duration} od \emph{object lifetime}:
\begin{itemize}
\item obiekt lokalny ma zazwyczaj automatyczny czas przechowywania;
\item obiekt z \texttt{static} ma statyczny czas przechowywania;
\item obiekt \texttt{thread\_local} ma czas przechowywania wątku;
\item pamięć otrzymana przez \texttt{new} lub alokator jest pamięcią
      o dynamicznym czasie przechowywania.
\end{itemize}
Nie należy utożsamiać tych kategorii z gwarantowanymi fizycznymi
obszarami „stos” i „sterta”: te nazwy opisują zwyczajowe implementacje,
a nie pełny model abstrakcyjny języka.

\begin{lstlisting}
#include <memory>

int global = 17;                 // static storage duration
thread_local int per_thread = 0; // thread storage duration

void demo() {
    int local = 5;               // automatic storage duration
    auto p = std::make_unique<int>(10);
    // Obiekt *p zyje w pamieci zarzadzanej dynamicznie.
}
\end{lstlisting}

\subsection{Reprezentacja obiektu i bezpieczne oglądanie bajtów}

C++ pozwala badać reprezentację bajtową obiektu przez
\texttt{char}, \texttt{unsigned char} lub \texttt{std::byte}.
Nie wolno jednak wyciągać z tego wniosku, że dowolny wskaźnik
\texttt{T*} można przekonwertować na \texttt{U*} i odczytać dane jako
\texttt{U}: takie postępowanie może naruszyć reguły wyrównania,
czasu życia i dostępu przez glvalue odpowiedniego typu.

\begin{lstlisting}
#include <array>
#include <cstddef>
#include <cstring>
#include <cstdint>

std::uint32_t x = 0x12345678;
std::array<std::byte, sizeof(x)> bytes{};
std::memcpy(bytes.data(), &x, sizeof(x));

std::uint32_t y = 0;
std::memcpy(&y, bytes.data(), sizeof(y));
// y == x
\end{lstlisting}

Kolejność bajtów zależy od \emph{endianness}. Można ją rozpoznać przez
\texttt{std::endian} (C++20, nagłówek \texttt{<bit>}), jednak do
serializacji najlepiej jawnie ustalać porządek bajtów, format pól
i wersję protokołu. Nie należy wysyłać całej struktury przez
\texttt{write(fd, \&object, sizeof(object))} jako przenośnego formatu:
padding, reprezentacja liczb i ABI mogą się różnić.

\subsection{\texttt{std::bit\_cast} i ograniczenia rzutowań}

\texttt{std::bit\_cast<To>(from)} (C++20, \texttt{<bit>}) tworzy wartość
\texttt{To} przez skopiowanie reprezentacji bitowej, o ile oba typy
są trywialnie kopiowalne i mają równy rozmiar. Dodatkowo reprezentacja
musi oznaczać poprawną wartość docelowego typu.

\begin{lstlisting}
#include <bit>
#include <cstdint>

float f = 1.0f;
// Przyklad zaklada sizeof(float) == sizeof(uint32_t).
// Na typowej platformie wynik to kod bitowy IEEE 754.
if constexpr (sizeof(float) == sizeof(std::uint32_t)) {
    auto bits = std::bit_cast<std::uint32_t>(f);
    (void)bits;
}
\end{lstlisting}

Nie mylmy \texttt{bit\_cast} z \texttt{reinterpret\_cast}:
pierwszy wytwarza \emph{wartość}, drugi może tylko zmienić sposób
interpretacji wskaźnika; sam nie tworzy obiektu docelowego.

\section{Layout struktury: \texttt{sizeof}, \texttt{alignof}, padding}

\begin{center}\small
\begin{tabularx}{\textwidth}{@{}p{0.33\textwidth}X@{}}
\toprule
Narzędzie & Co pokazuje lub zmienia \\
\midrule
\code{sizeof(T)} & Podaje rozmiar obiektu typu \code{T} w bajtach, wraz z ewentualnym paddingiem. \\
\code{alignof(T)} & Podaje wymagane wyrównanie adresu obiektu typu \code{T}. \\
\code{offsetof(T, field)} & Podaje przesunięcie składowej w typie o standardowym układzie. \\
\code{alignas(n)} & Żąda wyrównania co najmniej do granicy \code{n} dla obiektu lub typu. \\
\code{std::byte} & Reprezentuje jednostkę surowej pamięci bez znaczenia znaku tekstowego. \\
\bottomrule
\end{tabularx}
\end{center}
\label{sec:cpp10-layout}

\subsection{Skąd bierze się padding?}

Każdy typ obiektowy ma wymaganie wyrównania. Kompilator może wstawić
\emph{padding}, czyli nieużywane bajty, aby kolejne pola miały
odpowiednie adresy, a elementy tablicy struktur były poprawnie
wyrównane. Rozmieszczenie pól zależy od typu klasy i reguł języka;
konkretna liczba bajtów paddingu zależy od implementacji.

\begin{lstlisting}
#include <cstddef>
#include <iostream>

struct Unordered {
    char flag;
    double value;
    int count;
};

struct Ordered {
    double value;
    int count;
    char flag;
};

std::cout << sizeof(Unordered) << '\n';
std::cout << sizeof(Ordered) << '\n';
std::cout << alignof(Ordered) << '\n';
\end{lstlisting}

Na częstym ABI, przy \texttt{sizeof(double)==8},
\texttt{alignof(double)==8} i \texttt{sizeof(int)==4},
pierwszy typ zajmuje zwykle 24 bajty, a drugi 16. \textbf{Nie są to
wartości wymagane przez standard}. Samo przestawienie pól może
zmniejszyć zapotrzebowanie na pamięć milionów rekordów, ale może też
pogorszyć czytelność, stabilność ABI lub lokalność dostępu.

\begin{center}
\small
\begin{tabular}{|l|p{4.4cm}|p{4.4cm}|}
\hline
 & \textbf{Unordered (przykładowo)} & \textbf{Ordered (przykładowo)} \\
\hline
Początek & \texttt{flag}: 1 B & \texttt{value}: 8 B \\
\hline
Dalej & padding: 7 B & \texttt{count}: 4 B \\
\hline
Dalej & \texttt{value}: 8 B & \texttt{flag}: 1 B \\
\hline
Koniec & \texttt{count}: 4 B + padding: 4 B & padding: 3 B \\
\hline
Razem & 24 B & 16 B \\
\hline
\end{tabular}
\end{center}

\subsection{\texttt{offsetof}, standard-layout i ABI}

Dla typów \emph{standard-layout} makro \texttt{offsetof} pozwala
poznać przesunięcie pola od początku obiektu. Używanie go z innymi
typami klas od C++17 jest warunkowo wspierane, dlatego do analiz
przenośnych przyjmujemy tutaj standard-layout.

\begin{lstlisting}
#include <cstddef>
#include <cstdint>
#include <type_traits>

struct Header {
    std::uint32_t id;
    std::uint16_t version;
};
static_assert(std::is_standard_layout_v<Header>);
static_assert(offsetof(Header, version) >=
              sizeof(std::uint32_t));
\end{lstlisting}

W analizie ABI trzeba pamiętać o polach wirtualnych, dziedziczeniu,
przestrzeniach adresowych i konwencjach wywołania. Zmiana kolejności
pól publicznej struktury używanej w bibliotekach binarnych może
zmienić jej ABI nawet wtedy, gdy program źródłowy nadal się kompiluje.

\subsection{\texttt{alignof} i \texttt{alignas}}

\texttt{alignof(T)} zwraca wyrównanie typu \texttt{T};
\texttt{alignas(N)} pozwala narzucić silniejsze wyrównanie.
Należy zachować ograniczenia języka dotyczące dopuszczalnych wartości
\texttt{N}. Wyrównanie może zwiększyć rozmiar obiektu, co jest
istotne przy tablicach.

\begin{lstlisting}
#include <atomic>
#include <cstdint>

struct alignas(64) AlignedCounter {
    std::atomic<std::uint64_t> value{0};
};
static_assert(alignof(AlignedCounter) >= 64);
\end{lstlisting}

Wyrównanie do 64 bajtów ma sens tylko wtedy, gdy wynika z wymagań
sprzętu lub pomiarów. Nie jest to uniwersalna wielkość linii cache.
Ponadto \texttt{alignas} nie czyni dostępu do obiektu atomowym,
nie zastępuje \texttt{std::atomic} i samo nie zapewnia synchronizacji.

\subsection{Czy warto „pakować” struktury?}

Atrybuty kompilatora w rodzaju \texttt{\_\_attribute\_\_((packed))}
lub \texttt{\#pragma pack} nie stanowią przenośnych mechanizmów
standardowego C++20. Zmniejszenie rozmiaru może wymusić niewyrównany
dostęp, obniżyć wydajność, a na niektórych architekturach być
problemem poprawności. W kodzie protokołów lepiej często dekodować
pola jawnie z bufora bajtów i nie wiązać formatu z layoutem klasy.

\section{Wyrównane alokacje i czas życia obiektów}
\label{sec:cpp10-lifetime}

\subsection{Różnica między alokacją a konstruowaniem}

Klasyczne \texttt{new T(args)} wykonuje logicznie dwa zadania:
pozyskuje odpowiednio wyrównaną pamięć i inicjalizuje w niej obiekt.
Analogicznie \texttt{delete} niszczy obiekt i zwalnia pamięć.
Dla narzędzi niskopoziomowych warto rozdzielić te operacje.

\begin{lstlisting}
#include <new>

struct Widget { explicit Widget(int) {} };

void* mem = ::operator new(sizeof(Widget));
Widget* w = nullptr;
try {
    w = ::new (mem) Widget(5); // placement new
} catch (...) {
    ::operator delete(mem);
    throw;
}
w->~Widget();                 // koniec zycia obiektu
::operator delete(mem);      // zwolnienie samej pamieci
\end{lstlisting}

Taka konstrukcja jest poprawna dla \texttt{Widget} o wyrównaniu
obsługiwanym przez zwykłe \texttt{operator new}. Dla typów
\emph{over-aligned} należy korzystać z odpowiedniej alokacji
obsługującej \texttt{std::align\_val\_t} lub ze standardowego
alokatora, który uwzględnia wymagania wyrównania.

\paragraph{Dlaczego to nie jest zalecany idiom aplikacyjny?}
Obsługa wyjątków, wielokrotnych destruktorów, rozmiaru tablicy
i zgodności alokacji z dealokacją łatwo prowadzi do błędów.
W większości kodu wystarczają \texttt{std::vector},
\texttt{std::make\_unique} i obiekty automatyczne. Rozdzielanie
alokacji od konstruowania ma sens głównie przy implementowaniu
kontenerów, pul pamięci i innych abstrakcji zasobów.

\subsection{\texttt{std::construct\_at} i \texttt{std::destroy\_at} (C++20)}

Od C++20 można użyć funkcji z nagłówka \texttt{<memory>} do
jawnego rozpoczęcia i zakończenia czasu życia obiektu w odpowiednio
przygotowanej pamięci.

\begin{lstlisting}
#include <cstddef>
#include <memory>

struct Tracked {
    int value;
    explicit Tracked(int n) : value(n) {}
};

alignas(Tracked) std::byte buffer[sizeof(Tracked)];
Tracked* obj = std::construct_at(
    reinterpret_cast<Tracked*>(buffer), 42);
// ... korzystamy z *obj ...
std::destroy_at(obj);
\end{lstlisting}

\textbf{Ważne:} \texttt{std::destroy\_at} kończy czas życia obiektu, ale
nie zwalnia jego pamięci. Nie wolno wykonywać dwóch destrukcji tego
samego żywego obiektu ani odwoływać się do obiektu po destrukcji.
Ponadto rozmiar i wyrównanie bufora muszą odpowiadać konstruowanemu
typowi; sam \texttt{reinterpret\_cast} nie naprawia niewłaściwego
wyrównania.

\subsection{Wyrównany \texttt{operator new}}

Dla typów o szczególnie silnym wyrównaniu istnieją przeciążenia
alokacji i dealokacji wykorzystujące \texttt{std::align\_val\_t}
(C++17). Zwykłe \texttt{new Overaligned{}} wybiera odpowiedni
mechanizm automatycznie.

\begin{lstlisting}
#include <new>

struct alignas(128) Block { char data[128]; };
Block* b = new Block{};
delete b;
\end{lstlisting}

Jeśli ręcznie używamy \texttt{::operator new} z argumentem
\texttt{std::align\_val\_t\{128\}}, musimy dobrać zgodną operację
\texttt{::operator delete} z parametrem wyrównania. Mieszanie tych
rodzin funkcji jest błędem.

\subsection{Ponowne użycie tej samej pamięci i \texttt{std::launder}}

W tej samej pamięci można zakończyć życie jednego obiektu, a potem
rozpocząć życie kolejnego, ale dostęp przez wcześniej utworzone
wskaźniki podlega szczegółowym regułom języka. \texttt{std::launder}
(C++17) ma zastosowanie w określonych sytuacjach, gdy nowy obiekt
nie jest przezroczystym zamiennikiem poprzedniego; nie jest
uniwersalnym remedium na aliasing, nieprawidłowy czas życia
czy źle wyrównaną pamięć.

\paragraph{Reguła praktyczna.}
W kodzie aplikacyjnym nie sięgaj po \texttt{std::launder} bez
konkretnego problemu z czasem życia i bez zrozumienia odpowiednich
reguł. Preferuj własność zasobu wyrażoną typem RAII.

\section{Alokatory: pamięć, obiekty i bezpieczeństwo wyjątkowe}
\label{sec:cpp10-allocators}

\subsection{Standardowy model alokatora}

Kontenery standardowe używają \emph{alokatorów} do uzyskiwania
pamięci na swoje elementy. Alokator nie jest tym samym co funkcja
\texttt{new} i nie musi budować obiektu w chwili rezerwacji pamięci.
\texttt{std::allocator\_traits} zapewnia jednolity interfejs do
alokowania, konstruowania i niszczenia obiektów.

\begin{lstlisting}
#include <memory>

struct Item { explicit Item(int v) : value(v) {} int value; };
std::allocator<Item> alloc;
Item* p = alloc.allocate(2); // pamiec na dwa obiekty

std::allocator_traits<decltype(alloc)>::construct(alloc, p, 5);
std::allocator_traits<decltype(alloc)>::construct(alloc, p + 1, 7);

// teraz dwa obiekty sa zywe
std::allocator_traits<decltype(alloc)>::destroy(alloc, p + 1);
std::allocator_traits<decltype(alloc)>::destroy(alloc, p);
alloc.deallocate(p, 2);
\end{lstlisting}

Przykład dla przejrzystości zakłada, że oba konstruktory nie rzucają
wyjątków. \textbf{Pełna implementacja kontenera} musi zliczać
pomyślnie skonstruowane obiekty i w razie wyjątku zniszczyć tylko je,
a potem zwolnić pamięć. Właśnie z tego powodu lepiej opierać się na
gotowych kontenerach niż ręcznie pisać alokację elementów.

\subsection{Rezerwacja pojemności i realokacja \texttt{vector}}

\texttt{std::vector} przechowuje elementy w pamięci ciągłej, dzięki
czemu zapewnia dobrą lokalność dostępu. \texttt{reserve} może
ograniczyć liczbę realokacji, jeśli znamy przewidywaną wielkość
kontenera.

\begin{lstlisting}
#include <vector>

std::vector<int> values;
values.reserve(10'000);
for (int i = 0; i < 10'000; ++i)
    values.push_back(i);
\end{lstlisting}

\texttt{reserve(n)} nie zmienia \texttt{size()}, lecz zapewnia
\texttt{capacity() >= n}, chyba że zgłoszony zostanie wyjątek.
Realokacja unieważnia wskaźniki, referencje i iteratory do elementów.
Nie jest dobrym pomysłem wywoływanie \texttt{reserve(size()+1)}
przy każdym dodaniu elementu: można w ten sposób utracić
korzystną politykę geometrycznego wzrostu pojemności.

\subsection{Dlaczego własny alokator to nie pierwszy krok?}

Własny alokator w szablonie kontenera wiąże się z wymaganiami
dotyczącymi propagowania przy kopiowaniu, przenoszeniu i zamianie
kontenerów oraz z kwestią porównywania alokatorów.
Dla większości zastosowań wygodniejsza jest warstwa
\texttt{std::pmr}, którą poznamy teraz.

\section{\texttt{std::pmr}: polimorficzne zasoby pamięci}

\begin{center}\small
\begin{tabularx}{\textwidth}{@{}p{0.37\textwidth}X@{}}
\toprule
Mechanizm & Zastosowanie i ograniczenie \\
\midrule
\code{pmr::memory_resource} & Interfejs strategii przydziału i zwalniania pamięci używany przez alokatory PMR. \\
\code{pmr::vector<T>} & Kontener, którego alokator korzysta z wybranego zasobu pamięci. \\
\code{monotonic_buffer_resource} & Przydziela kolejne bloki szybko; pojedyncze zwolnienia nie odzyskują miejsca, zasób zwalnia je zbiorczo. \\
\code{unsynchronized_pool_resource} & Pula małych bloków bez synchronizacji między wątkami. \\
\code{synchronized_pool_resource} & Pula z synchronizacją, odpowiednia do współdzielenia między wątkami. \\
\bottomrule
\end{tabularx}
\end{center}
\label{sec:cpp10-pmr}

\subsection{Alokator jako decyzja w czasie działania}

\texttt{std::pmr} (C++17, \texttt{<memory\_resource>}) oddziela
\emph{rodzaj kontenera} od \emph{sposobu przydzielania pamięci}.
\texttt{std::pmr::vector<T>} jest aliasem \texttt{std::vector<T, ...>}
z alokatorem polimorficznym. Zasób \texttt{memory\_resource}
można wybierać w czasie działania programu.

\begin{lstlisting}
#include <memory_resource>
#include <vector>

std::pmr::monotonic_buffer_resource arena;
std::pmr::vector<int> numbers{&arena};
for (int i = 0; i < 100; ++i)
    numbers.push_back(i);
\end{lstlisting}

Kontener przechowuje \emph{wskaźnik do zasobu}, ale nie przejmuje
własności zasobu pamięci. Zasób musi
pozostać żywy wystarczająco długo, aby obsłużyć wszystkie powiązane
alokacje i dealokacje. \textbf{Nie wolno niszczyć lub resetować
areny, gdy żywe obiekty wciąż korzystają z jej pamięci}.

\subsection{Arena monotoniczna: wiele alokacji, jedno zwolnienie}

\texttt{std::pmr::monotonic\_buffer\_resource} wydaje kolejne
fragmenty pamięci. Jej \texttt{deallocate} nie zwalnia pojedynczych
bloków; odzyskanie następuje przy \texttt{release()} lub destrukcji
zasobu. To pasuje do wielu obiektów budowanych razem i usuwanych
w jednej fazie, np. drzewa parsowania pojedynczego dokumentu.

\begin{lstlisting}
#include <cstddef>
#include <memory_resource>
#include <vector>

alignas(std::max_align_t) std::byte storage[4096];
std::pmr::monotonic_buffer_resource arena(
    storage, sizeof(storage),
    std::pmr::null_memory_resource());

{
    std::pmr::vector<int> data{&arena};
    for (int i = 0; i < 32; ++i) data.push_back(i);
} // kontener znika; jego deallocate nie uwalnia blokow
// arena jest teraz bezpieczna do zniszczenia lub release()
\end{lstlisting}

\texttt{null\_memory\_resource()} wymusza niepowodzenie alokacji,
gdy lokalny bufor się wyczerpie (rzucane jest \texttt{std::bad\_alloc}).
Jeżeli zamiast niego użyjemy zwykłego zasobu nadrzędnego
(\emph{upstream}), arena może dobierać kolejne bufory.

\paragraph{Ważna konsekwencja.}
Gdy \texttt{vector} zwiększa pojemność, porzuca stary blok.
W przypadku areny monotonicznej ten blok zwykle pozostaje zajęty
aż do końca całej fazy. Arena nie oznacza zatem automatycznie
najmniejszego zużycia pamięci; warto stosować \texttt{reserve}.

\subsection{Zasoby pulowe: \texttt{unsynchronized\_pool\_resource}}

Gdy obiekty są często alokowane i usuwane w różnych momentach,
przydatne bywają pule przechowujące bloki o podobnych rozmiarach.

\begin{lstlisting}
#include <memory_resource>
#include <vector>

std::pmr::unsynchronized_pool_resource pool;
std::pmr::vector<int> items{&pool};
items.resize(500);
\end{lstlisting}

\texttt{unsynchronized\_pool\_resource} nie zapewnia bezpieczeństwa
jednoczesnego użycia z wielu wątków. Do współdzielonego użycia
istnieje \texttt{synchronized\_pool\_resource}. Wybór zależy od
profilu alokacji i pomiarów: pule nie są uniwersalnie szybsze.

\subsection{Własny zasób: liczymy rzeczywiste alokacje}

Wariant diagnostyczny może zliczać żądania pamięci i delegować je
do \texttt{new\_delete\_resource()}. Trzy wirtualne funkcje
\texttt{do\_allocate}, \texttt{do\_deallocate}, \texttt{do\_is\_equal}
tworzą podstawę własnej implementacji.

\begin{lstlisting}
#include <cstddef>
#include <memory_resource>

struct CountingResource : std::pmr::memory_resource {
    std::pmr::memory_resource* upstream =
        std::pmr::new_delete_resource();
    std::size_t allocations = 0;
    std::size_t bytes = 0;
private:
    void* do_allocate(std::size_t n,
                      std::size_t alignment) override {
        void* p = upstream->allocate(n, alignment);
        ++allocations;
        bytes += n;
        return p;
    }
    void do_deallocate(void* p, std::size_t n,
                       std::size_t alignment) override {
        upstream->deallocate(p, n, alignment);
    }
    bool do_is_equal(
        const std::pmr::memory_resource& other
    ) const noexcept override {
        return this == &other;
    }
};
\end{lstlisting}

\begin{lstlisting}
CountingResource stats;
{
    std::pmr::vector<int> data{&stats};
    for (int i = 0; i < 100; ++i) data.push_back(i);
}
// stats.allocations i stats.bytes opisuja zadania
// przeslane do tego konkretnego zasobu.
\end{lstlisting}

\textbf{Uwaga:} nie są to wszystkie alokacje programu, a
\texttt{bytes} oznacza sumę żądanych rozmiarów, nie największe
jednoczesne zużycie pamięci. Warto liczyć także dealokacje,
\emph{peak usage} i rozkład rozmiarów bloków.

\subsection{Zagnieżdżone kontenery i propagacja zasobu}

\begin{lstlisting}
#include <memory_resource>
#include <string>
#include <vector>

CountingResource stats;
std::pmr::vector<std::pmr::string> names{&stats};
names.emplace_back("alpha");
names.emplace_back("a rather long text that needs allocation");
\end{lstlisting}

\texttt{pmr::string} również korzysta z polimorficznego alokatora,
ale dokładne zachowanie małych napisów zależy od implementacji
(\emph{small string optimization}). Nie należy zakładać, że każdy
\texttt{emplace\_back} przydziela pamięć ani że krótki napis wymaga
alokacji dynamicznej.

\section{Lokalność danych i hierarchia cache}
\label{sec:cpp10-cache}

\subsection{Dlaczego asymptotyka nie wystarcza?}

Dwa algorytmy o złożoności \(O(n)\) mogą mieć bardzo różną wydajność,
jeżeli pierwszy czyta sąsiadujące elementy tablicy, a drugi podąża
za wskaźnikami do rozproszonych obiektów. Procesor pobiera dane
porcjami (\emph{cache lines}); dostęp sekwencyjny sprzyja
prefetchingowi, zaś losowy dostęp bywa ograniczony opóźnieniem pamięci.

Nie ma jednego rozmiaru linii cache dla wszystkich procesorów ani
jednej wartości kosztu trafienia i chybienia. Zależą one od sprzętu,
aktualnego obciążenia i rodzaju dostępu.

\begin{lstlisting}
#include <numeric>
#include <vector>

std::vector<int> v(100000, 1);
long long sum = std::accumulate(v.begin(), v.end(), 0LL);
// Dostep sekwencyjny do pamieci ciaglej.
\end{lstlisting}

Wskaźniki do osobno alokowanych małych węzłów mogą pogorszyć
lokalność i powiększyć narzut pamięciowy. Nie oznacza to,
że \texttt{std::vector} jest zawsze lepszy od \texttt{std::list}:
wybór zależy od wymaganych operacji, unieważniania iteratorów
i rzeczywistego profilu obciążenia.

\subsection{AoS kontra SoA}

W układzie \emph{Array of Structures} (AoS) przechowujemy obok siebie
wszystkie składowe kolejnych rekordów. W \emph{Structure of Arrays}
(SoA) przechowujemy osobne tablice poszczególnych składowych.

\begin{lstlisting}
#include <vector>

struct PointAoS { float x, y, z, intensity; };
std::vector<PointAoS> aos;

struct PointSoA {
    std::vector<float> x, y, z, intensity;
};
PointSoA soa;
\end{lstlisting}

Gdy obliczamy wyłącznie sumę współrzędnych \texttt{x}, układ SoA
umożliwia odczytywanie samych potrzebnych danych, a nie całych
rekordów. Ułatwia to też wektoryzację. Jeśli natomiast w każdej
operacji potrzebujemy wszystkich pól jednego rekordu, AoS może być
korzystniejszy. Wydajność zależy od rozmiaru danych, kodu i sprzętu.

\begin{center}
\small
\begin{tabular}{|p{6.0cm}|p{6.0cm}|}
\hline
\textbf{AoS: rekordy obok siebie} & \textbf{SoA: kolumny obok siebie} \\
\hline
\texttt{[x1 y1 z1 i1] [x2 y2 z2 i2]} &
\texttt{[x1 x2 ...] [y1 y2 ...]} \\
\hline
Wygodny dostęp do całego punktu & Wygodny dostęp do jednej cechy wszystkich punktów \\
\hline
Często mniej kontenerów i prostszy model & Często lepsze wykorzystanie cache w obliczeniach kolumnowych \\
\hline
\end{tabular}
\end{center}

\subsection{Stride, skoki i przewidywanie dostępu}

Przechodzenie po pamięci ze skokiem \(k\) oznacza, że kolejne
operacje mogą dotykać odległych linii cache. Dla macierzy w układzie
wierszowym zwykle korzystniej jest przechodzić po kolumnach
\emph{wewnątrz pojedynczego wiersza}, a następnie po wierszach.
Zanim zmienimy kolejność pętli, należy uwzględnić zależności
obliczeniowe, autowektoryzację i to, czy kompilator już dokonał
danej transformacji.

\begin{lstlisting}
#include <vector>

constexpr std::size_t rows = 100;
constexpr std::size_t cols = 200;
std::vector<double> matrix(rows * cols, 1.0);
double sum = 0.0;
for (std::size_t i = 0; i < rows; ++i)
    for (std::size_t j = 0; j < cols; ++j)
        sum += matrix[i * cols + j];
\end{lstlisting}

\section{False sharing i współbieżne liczniki}
\label{sec:cpp10-false-sharing}

\subsection{Na czym polega false sharing?}

Jeżeli dwa wątki wielokrotnie zapisują \emph{różne} zmienne, ale
zmienne znajdują się w tej samej linii cache, protokół spójności
pamięci podręcznej może wymuszać kosztowną wymianę własności tej
linii między rdzeniami. Nazywamy to \emph{false sharing}.
Nie oznacza to \emph{data race}, jeżeli poszczególne zmienne są
atomowe lub inaczej poprawnie synchronizowane.

\begin{lstlisting}
#include <atomic>
#include <cstdint>

struct NaiveCounters {
    std::atomic<std::uint64_t> first{0};
    std::atomic<std::uint64_t> second{0};
};
\end{lstlisting}

Dwa liczniki mogą być blisko siebie, jednak dokładny layout obiektu
i mapowanie do linii cache nie są gwarantowane. Na niektórych
platformach ich równoczesna modyfikacja będzie kosztowna.

\subsection{Rozdzielenie liczników w pamięci}

\begin{lstlisting}
#include <atomic>
#include <cstdint>

struct alignas(64) IsolatedAtomic {
    std::atomic<std::uint64_t> value{0};
};
struct SeparatedCounters {
    IsolatedAtomic first, second;
};
\end{lstlisting}

Tu użyliśmy 64 bajtów wyłącznie jako \emph{przykładowej} wartości.
Od C++17 w nagłówku \texttt{<new>} istnieją dwie stałe
opisujące zalecane odstępy dla danych często modyfikowanych przez
różne wątki (\emph{destructive interference}) oraz dla danych,
które powinny pozostawać blisko siebie (\emph{constructive
interference}). Pełną nazwę pierwszej z nich pokazuje kod poniżej.
Wartości stałych zależą od implementacji; nie stanowią pomiaru
rozmiaru linii cache w uruchomionym procesorze.

\begin{lstlisting}
#include <atomic>
#include <new>

struct alignas(
    std::hardware_destructive_interference_size
) Counter {
    std::atomic<unsigned> value{0};
};
\end{lstlisting}

Rozdzielenie liczników zwiększa rozmiar struktury. Przy milionie
liczników koszt pamięciowy może przewyższyć korzyść. Technika ma
największy sens, gdy oddzielne wątki często zapisują niewielką
liczbę niezależnych pól.

\subsection{False sharing a semantyka atomowa}

\texttt{std::memory\_order\_relaxed} wystarczy do niektórych
niezależnych liczników, gdy zależy nam wyłącznie na atomowym
zliczaniu; nie gwarantuje jednak synchronizacji innych danych.
Usunięcie false sharing nie zastępuje poprawnego modelu pamięci
współbieżnej. Szczegółowe zasady operacji atomowych były omawiane
w rozdziale~1 o współbieżności.

\section{Jak rzetelnie mierzyć wydajność?}
\label{sec:cpp10-benchmark}

\subsection{Zasady pomiarów}

Najważniejsze pytanie brzmi nie „czy ten kod wygląda szybciej?”,
lecz „czy w rzeczywistym scenariuszu uzyskujemy powtarzalną poprawę?”.
Doświadczenie powinno uwzględniać:
\begin{enumerate}
\item kompilację zoptymalizowaną (np. \texttt{-O2} lub \texttt{-O3});
\item \emph{warm-up} i wielokrotne uruchomienia;
\item jednakowe dane oraz taki sam rezultat dla obu wariantów;
\item ochronę przed usunięciem testowanej pracy przez optymalizator;
\item pomiar rozkładu wyników, a nie pojedynczego odczytu;
\item ograniczenie wpływu alokacji, I/O, schedulerów i skalowania CPU;
\item kontrolę zużycia pamięci, a nie tylko czasu;
\item narzędzia profilujące (np. Linux \texttt{perf}, odpowiedniki
      sprzętowe lub profilery danego środowiska).
\end{enumerate}

\subsection{Minimalna funkcja pomiarowa}

\begin{lstlisting}
#include <chrono>
#include <utility>

template<class F>
auto measure(F&& f) {
    auto start = std::chrono::steady_clock::now();
    auto result = std::forward<F>(f)();
    auto stop = std::chrono::steady_clock::now();
    return std::pair{result, stop - start};
}
\end{lstlisting}

Taka funkcja przydaje się do nauki, ale \emph{nie jest kompletna}
biblioteka benchmarkowa. Dla zbyt krótkich fragmentów narzut
pomiaru może przeważać nad czasem pracy. Dla liczb zmiennoprzecinkowych
porównywanie różnych kolejności sumowania wymaga ostrożności,
ponieważ zmiana kolejności działań może zmienić zaokrąglenia.

\subsection{Sanitizatory i wydajność}

Sanitizatory, np. AddressSanitizer i UndefinedBehaviorSanitizer,
są narzędziami do badania poprawności pamięci i innych błędów.
Wprowadzają narzut, dlatego \emph{pomiarów porównawczych wydajności}
nie należy bezrefleksyjnie wykonywać na buildach sanitizowanych.
Bezpieczny cykl pracy wygląda następująco: test poprawności,
profilowanie zwykłego buildu, zmiana jednej rzeczy, ponowny test
poprawności i ponowny pomiar.

\section{Przykładowy projekt: szybkie przetwarzanie wielu punktów}
\label{sec:cpp10-project}

Przyjmijmy, że program przechowuje milion punktów \((x,y,z,I)\)
i dla każdego oblicza tylko prosty filtr zależny od \(x\) i \(I\).
Możemy rozważyć dwa modele:
\begin{enumerate}
\item \textbf{AoS:} \texttt{vector<Point>} --- prostszy interfejs,
      wygodny dostęp do pełnego punktu.
\item \textbf{SoA:} osobne \texttt{vector<float>} na \(x,y,z,I\) ---
      możliwy mniejszy transfer danych potrzebnych do filtru.
\end{enumerate}

Następnie oceniamy, czy warto użyć \texttt{reserve} przed zapełnieniem
kontenerów; jeśli rekordy są tworzone na potrzeby jednego zapytania,
możemy sprawdzić arenę monotoniczną. Trzeba jednak rozróżnić
\emph{koszt wypełnienia struktur} od \emph{kosztu przetwarzania}: arena
zmienia zachowanie alokacji, nie gwarantuje poprawy pracy procesora
nad danymi.

Dla kodu numerycznego duży wpływ może mieć też wybór kolejności
wymiarów macierzy, reprezentacji sąsiedztwa w siatce czy sposobu
obliczania operatora na konkretnym punkcie zamiast konstruowania
dużych obiektów pośrednich. Właściwa reprezentacja danych potrafi
dać większy zysk niż drobne optymalizacje pojedynczych instrukcji.

\section{Jak dobrać technikę?}

\begin{center}
\small
\begin{tabular}{|p{5.7cm}|p{6.7cm}|}
\hline
\textbf{Problem} & \textbf{Pierwszy sensowny krok} \\
\hline
Za dużo dynamicznych alokacji & \texttt{reserve}, ograniczenie obiektów pośrednich, pomiar alokacji \\
\hline
Wiele obiektów znika razem & \texttt{pmr::monotonic\_buffer\_resource} po kontroli czasu życia \\
\hline
Dużo małych alokacji o różnej żywotności & \texttt{pmr::...pool\_resource} i pomiar \\
\hline
Praca nad jedną kolumną wielu rekordów & Porównanie AoS z SoA \\
\hline
Losowe skoki po pamięci & Sprawdzenie lokalności, zmian formatu danych i blokowania pętli \\
\hline
Niezależne atomiki spowalniają się nawzajem & Analiza false sharing i separacja intensywnie zapisywanych pól \\
\hline
Trzeba uzyskać silne wyrównanie & \texttt{alignas}, odpowiednia alokacja i kontrola \texttt{alignof} \\
\hline
Trzeba napisać kontener lub pulę & \texttt{allocator\_traits}, zarządzanie czasem życia, RAII \\
\hline
Potrzeba przenośnej serializacji & Jawny format bajtowy, bez kopiowania całych struktur \\
\hline
\end{tabular}
\end{center}

\section{Typowe błędy i pułapki}

\begin{enumerate}
\item \textbf{Wniosek „\texttt{sizeof(T)} to suma pól”.}
      Padding może zwiększyć rozmiar, a reguły klas mogą być złożone.
\item \textbf{Uznanie layoutu z jednej kompilacji za standard.}
      ABI, wyrównanie i rozmiary typów różnią się między platformami.
\item \textbf{Nieprawidłowe \texttt{reinterpret\_cast}.}
      Sama zmiana typu wskaźnika nie legalizuje dostępu do danych.
\item \textbf{Mylenie pamięci z czasem życia obiektu.}
      \texttt{allocate} nie zastępuje konstruowania obiektów klasowych.
\item \textbf{Zwalnianie pamięci bez destruktora.}
      Dla nietrywialnych typów może wyciekać zasób zarządzany przez obiekt.
\item \textbf{Niezgodna para alokacja--dealokacja.}
      Szczególnie ważne przy over-alignment i ręcznej alokacji.
\item \textbf{Pozostawienie kontenera po zniszczeniu areny \texttt{pmr}.}
      Jego pamięć staje się niepoprawna; wynik to niezdefiniowane zachowanie.
\item \textbf{Wywołanie \texttt{release()} przy żywych obiektach.}
      Arena może unieważnić ich pamięć bez wywoływania destruktorów.
\item \textbf{Założenie, że \texttt{monotonic\_buffer\_resource} oszczędza
      pamięć.} Zwolnienie pojedynczych bloków zwykle nie następuje.
\item \textbf{Współdzielenie \texttt{unsynchronized\_pool\_resource}
      między wątkami bez ochrony.} Zasób nie jest thread-safe.
\item \textbf{Optymalizacja poprzez mniejszą strukturę bez pomiaru.}
      Mniejszy \texttt{sizeof} nie gwarantuje krótszego czasu wykonania.
\item \textbf{Rzekoma optymalizacja przez \texttt{volatile}.}
      \texttt{volatile} nie zapewnia synchronizacji między wątkami.
\item \textbf{Pomiar jednego przebiegu.} Scheduler, częstotliwości
      zegara i cache mogą silnie zniekształcić wynik.
\item \textbf{Użycie sanitizatora jako miary szybkości.}
      Narzut instrumentacji zmienia charakterystykę programu.
\end{enumerate}

\paragraph{Pytania kontrolne.}
\begin{enumerate}
\item Czym różnią się \texttt{sizeof(T)} i \texttt{alignof(T)}?
\item Dlaczego zmiana kolejności pól może zmienić rozmiar struktury?
\item Czy placement \texttt{new} przydziela pamięć?
\item Jaka jest różnica między \texttt{destroy\_at} a \texttt{deallocate}?
\item Kiedy arena monotoniczna może zużywać więcej pamięci niż oczekujemy?
\item Dlaczego \texttt{pmr::vector<T>} nie jest właścicielem zasobu?
\item Kiedy SoA może przyspieszyć obliczenia na rekordach?
\item Czy false sharing i data race są tym samym problemem?
\item Dlaczego \texttt{std::memory\_order\_relaxed} nie rozwiązuje
      wszystkich problemów synchronizacji?
\item Jak odróżnić realną poprawę wydajności od szumu pomiarowego?
\end{enumerate}

\paragraph{Dalsza lektura.}
Dokumentacja \emph{cppreference} obejmuje hasła \emph{Object},
\emph{Object lifetime}, \texttt{alignof}, \texttt{alignas},
\texttt{std::construct\_at} oraz \texttt{allocator\_traits}.
Osobno warto przeczytać opis zasobów \texttt{std::pmr}
i własności \emph{hardware interference size}.
W praktyce polecam analizę śladów alokacji, profilery
wydajności i \emph{Compiler Explorer} do porównywania
wygenerowanego kodu maszynowego.


\chapter[Organizacja pamięci i wydajność]{Organizacja pamięci, wyrównanie i wydajność niskopoziomowa w C++}
\label{chap:cpp-memory-performance}

\section{Po co analizować organizację pamięci?}

W rozdziałach o wskaźnikach (rozdział~3), semantyce przenoszenia
(rozdział~6) i iteratorach (rozdział~8) interesowało nas głównie
\emph{poprawne zarządzanie obiektami}. Teraz przyjrzymy się temu, jak
obiekty są reprezentowane w pamięci oraz w jaki sposób wybory
programisty wpływają na liczbę alokacji i lokalność dostępu do danych.

Rozróżniajmy trzy poziomy opisu:
\begin{enumerate}
\item \textbf{Gwarancje języka C++:} rozmiar obiektu, wymagane wyrównanie,
      czas życia obiektu, legalność dostępu i zasady alokowania pamięci.
\item \textbf{Szczegóły implementacji:} rozmieszczenie pól, rozmiary
      wskaźników, zachowanie konkretnego alokatora czy układ binarny ABI.
\item \textbf{Właściwości sprzętu:} linie cache, hierarchia pamięci,
      prefetching, przepustowość i zjawisko \emph{false sharing}.
\end{enumerate}

Kod poprawny według standardu może być wolny; szybki kod działający
na jednym komputerze może być powolny na innym. Dlatego optymalizacja
wymaga najpierw \emph{pomiaru}, a następnie weryfikacji hipotezy.

\section{Obiekt, bajt, adres i czas życia}
\label{sec:cpp10-object}

\subsection{Pamięć nie jest tym samym co obiekt}

W modelu C++ \emph{obszar pamięci} może istnieć bez żywego obiektu typu
\texttt{T}. Obiekt charakteryzują między innymi typ, rozmiar, czas życia
i, dla pewnych klas, odpowiednio wywołane konstruktory i destruktory.
Samo uzyskanie adresu wystarczająco dużego bloku nie oznacza, że
wolno traktować ten blok jak dowolny obiekt C++.

\begin{lstlisting}
#include <cstddef>
#include <type_traits>

static_assert(sizeof(char) == 1);
static_assert(sizeof(std::byte) == 1);
static_assert(sizeof(int) >= 1);

struct Point { double x, y; };
static_assert(sizeof(Point) >= 2 * sizeof(double));
\end{lstlisting}

\texttt{sizeof(T)} podaje rozmiar obiektu w \emph{bajtach C++}.
Jeden bajt nie musi mieć ośmiu bitów; ich liczbę określa
\texttt{CHAR\_BIT} z nagłówka \texttt{<climits>}.
Rozmiar typu obiektowego jest dodatni. Każdy obiekt ma też
wymagane \emph{wyrównanie}, które omówimy dalej.

\paragraph{Czas przechowywania a czas życia.}
Należy odróżnić \emph{storage duration} od \emph{object lifetime}:
\begin{itemize}
\item obiekt lokalny ma zazwyczaj automatyczny czas przechowywania;
\item obiekt z \texttt{static} ma statyczny czas przechowywania;
\item obiekt \texttt{thread\_local} ma czas przechowywania wątku;
\item pamięć otrzymana przez \texttt{new} lub alokator jest pamięcią
      o dynamicznym czasie przechowywania.
\end{itemize}
Nie należy utożsamiać tych kategorii z gwarantowanymi fizycznymi
obszarami „stos” i „sterta”: te nazwy opisują zwyczajowe implementacje,
a nie pełny model abstrakcyjny języka.

\begin{lstlisting}
#include <memory>

int global = 17;                 // static storage duration
thread_local int per_thread = 0; // thread storage duration

void demo() {
    int local = 5;               // automatic storage duration
    auto p = std::make_unique<int>(10);
    // Obiekt *p zyje w pamieci zarzadzanej dynamicznie.
}
\end{lstlisting}

\subsection{Reprezentacja obiektu i bezpieczne oglądanie bajtów}

C++ pozwala badać reprezentację bajtową obiektu przez
\texttt{char}, \texttt{unsigned char} lub \texttt{std::byte}.
Nie wolno jednak wyciągać z tego wniosku, że dowolny wskaźnik
\texttt{T*} można przekonwertować na \texttt{U*} i odczytać dane jako
\texttt{U}: takie postępowanie może naruszyć reguły wyrównania,
czasu życia i dostępu przez glvalue odpowiedniego typu.

\begin{lstlisting}
#include <array>
#include <cstddef>
#include <cstring>
#include <cstdint>

std::uint32_t x = 0x12345678;
std::array<std::byte, sizeof(x)> bytes{};
std::memcpy(bytes.data(), &x, sizeof(x));

std::uint32_t y = 0;
std::memcpy(&y, bytes.data(), sizeof(y));
// y == x
\end{lstlisting}

Kolejność bajtów zależy od \emph{endianness}. Można ją rozpoznać przez
\texttt{std::endian} (C++20, nagłówek \texttt{<bit>}), jednak do
serializacji najlepiej jawnie ustalać porządek bajtów, format pól
i wersję protokołu. Nie należy wysyłać całej struktury przez
\texttt{write(fd, \&object, sizeof(object))} jako przenośnego formatu:
padding, reprezentacja liczb i ABI mogą się różnić.

\subsection{\texttt{std::bit\_cast} i ograniczenia rzutowań}

\texttt{std::bit\_cast<To>(from)} (C++20, \texttt{<bit>}) tworzy wartość
\texttt{To} przez skopiowanie reprezentacji bitowej, o ile oba typy
są trywialnie kopiowalne i mają równy rozmiar. Dodatkowo reprezentacja
musi oznaczać poprawną wartość docelowego typu.

\begin{lstlisting}
#include <bit>
#include <cstdint>

float f = 1.0f;
// Przyklad zaklada sizeof(float) == sizeof(uint32_t).
// Na typowej platformie wynik to kod bitowy IEEE 754.
if constexpr (sizeof(float) == sizeof(std::uint32_t)) {
    auto bits = std::bit_cast<std::uint32_t>(f);
    (void)bits;
}
\end{lstlisting}

Nie mylmy \texttt{bit\_cast} z \texttt{reinterpret\_cast}:
pierwszy wytwarza \emph{wartość}, drugi może tylko zmienić sposób
interpretacji wskaźnika; sam nie tworzy obiektu docelowego.

\section{Layout struktury: \texttt{sizeof}, \texttt{alignof}, padding}

\begin{center}\small
\begin{tabularx}{\textwidth}{@{}p{0.33\textwidth}X@{}}
\toprule
Narzędzie & Co pokazuje lub zmienia \\
\midrule
\code{sizeof(T)} & Podaje rozmiar obiektu typu \code{T} w bajtach, wraz z ewentualnym paddingiem. \\
\code{alignof(T)} & Podaje wymagane wyrównanie adresu obiektu typu \code{T}. \\
\code{offsetof(T, field)} & Podaje przesunięcie składowej w typie o standardowym układzie. \\
\code{alignas(n)} & Żąda wyrównania co najmniej do granicy \code{n} dla obiektu lub typu. \\
\code{std::byte} & Reprezentuje jednostkę surowej pamięci bez znaczenia znaku tekstowego. \\
\bottomrule
\end{tabularx}
\end{center}
\label{sec:cpp10-layout}

\subsection{Skąd bierze się padding?}

Każdy typ obiektowy ma wymaganie wyrównania. Kompilator może wstawić
\emph{padding}, czyli nieużywane bajty, aby kolejne pola miały
odpowiednie adresy, a elementy tablicy struktur były poprawnie
wyrównane. Rozmieszczenie pól zależy od typu klasy i reguł języka;
konkretna liczba bajtów paddingu zależy od implementacji.

\begin{lstlisting}
#include <cstddef>
#include <iostream>

struct Unordered {
    char flag;
    double value;
    int count;
};

struct Ordered {
    double value;
    int count;
    char flag;
};

std::cout << sizeof(Unordered) << '\n';
std::cout << sizeof(Ordered) << '\n';
std::cout << alignof(Ordered) << '\n';
\end{lstlisting}

Na częstym ABI, przy \texttt{sizeof(double)==8},
\texttt{alignof(double)==8} i \texttt{sizeof(int)==4},
pierwszy typ zajmuje zwykle 24 bajty, a drugi 16. \textbf{Nie są to
wartości wymagane przez standard}. Samo przestawienie pól może
zmniejszyć zapotrzebowanie na pamięć milionów rekordów, ale może też
pogorszyć czytelność, stabilność ABI lub lokalność dostępu.

\begin{center}
\small
\begin{tabular}{|l|p{4.4cm}|p{4.4cm}|}
\hline
 & \textbf{Unordered (przykładowo)} & \textbf{Ordered (przykładowo)} \\
\hline
Początek & \texttt{flag}: 1 B & \texttt{value}: 8 B \\
\hline
Dalej & padding: 7 B & \texttt{count}: 4 B \\
\hline
Dalej & \texttt{value}: 8 B & \texttt{flag}: 1 B \\
\hline
Koniec & \texttt{count}: 4 B + padding: 4 B & padding: 3 B \\
\hline
Razem & 24 B & 16 B \\
\hline
\end{tabular}
\end{center}

\subsection{\texttt{offsetof}, standard-layout i ABI}

Dla typów \emph{standard-layout} makro \texttt{offsetof} pozwala
poznać przesunięcie pola od początku obiektu. Używanie go z innymi
typami klas od C++17 jest warunkowo wspierane, dlatego do analiz
przenośnych przyjmujemy tutaj standard-layout.

\begin{lstlisting}
#include <cstddef>
#include <cstdint>
#include <type_traits>

struct Header {
    std::uint32_t id;
    std::uint16_t version;
};
static_assert(std::is_standard_layout_v<Header>);
static_assert(offsetof(Header, version) >=
              sizeof(std::uint32_t));
\end{lstlisting}

W analizie ABI trzeba pamiętać o polach wirtualnych, dziedziczeniu,
przestrzeniach adresowych i konwencjach wywołania. Zmiana kolejności
pól publicznej struktury używanej w bibliotekach binarnych może
zmienić jej ABI nawet wtedy, gdy program źródłowy nadal się kompiluje.

\subsection{\texttt{alignof} i \texttt{alignas}}

\texttt{alignof(T)} zwraca wyrównanie typu \texttt{T};
\texttt{alignas(N)} pozwala narzucić silniejsze wyrównanie.
Należy zachować ograniczenia języka dotyczące dopuszczalnych wartości
\texttt{N}. Wyrównanie może zwiększyć rozmiar obiektu, co jest
istotne przy tablicach.

\begin{lstlisting}
#include <atomic>
#include <cstdint>

struct alignas(64) AlignedCounter {
    std::atomic<std::uint64_t> value{0};
};
static_assert(alignof(AlignedCounter) >= 64);
\end{lstlisting}

Wyrównanie do 64 bajtów ma sens tylko wtedy, gdy wynika z wymagań
sprzętu lub pomiarów. Nie jest to uniwersalna wielkość linii cache.
Ponadto \texttt{alignas} nie czyni dostępu do obiektu atomowym,
nie zastępuje \texttt{std::atomic} i samo nie zapewnia synchronizacji.

\subsection{Czy warto „pakować” struktury?}

Atrybuty kompilatora w rodzaju \texttt{\_\_attribute\_\_((packed))}
lub \texttt{\#pragma pack} nie stanowią przenośnych mechanizmów
standardowego C++20. Zmniejszenie rozmiaru może wymusić niewyrównany
dostęp, obniżyć wydajność, a na niektórych architekturach być
problemem poprawności. W kodzie protokołów lepiej często dekodować
pola jawnie z bufora bajtów i nie wiązać formatu z layoutem klasy.

\section{Wyrównane alokacje i czas życia obiektów}
\label{sec:cpp10-lifetime}

\subsection{Różnica między alokacją a konstruowaniem}

Klasyczne \texttt{new T(args)} wykonuje logicznie dwa zadania:
pozyskuje odpowiednio wyrównaną pamięć i inicjalizuje w niej obiekt.
Analogicznie \texttt{delete} niszczy obiekt i zwalnia pamięć.
Dla narzędzi niskopoziomowych warto rozdzielić te operacje.

\begin{lstlisting}
#include <new>

struct Widget { explicit Widget(int) {} };

void* mem = ::operator new(sizeof(Widget));
Widget* w = nullptr;
try {
    w = ::new (mem) Widget(5); // placement new
} catch (...) {
    ::operator delete(mem);
    throw;
}
w->~Widget();                 // koniec zycia obiektu
::operator delete(mem);      // zwolnienie samej pamieci
\end{lstlisting}

Taka konstrukcja jest poprawna dla \texttt{Widget} o wyrównaniu
obsługiwanym przez zwykłe \texttt{operator new}. Dla typów
\emph{over-aligned} należy korzystać z odpowiedniej alokacji
obsługującej \texttt{std::align\_val\_t} lub ze standardowego
alokatora, który uwzględnia wymagania wyrównania.

\paragraph{Dlaczego to nie jest zalecany idiom aplikacyjny?}
Obsługa wyjątków, wielokrotnych destruktorów, rozmiaru tablicy
i zgodności alokacji z dealokacją łatwo prowadzi do błędów.
W większości kodu wystarczają \texttt{std::vector},
\texttt{std::make\_unique} i obiekty automatyczne. Rozdzielanie
alokacji od konstruowania ma sens głównie przy implementowaniu
kontenerów, pul pamięci i innych abstrakcji zasobów.

\subsection{\texttt{std::construct\_at} i \texttt{std::destroy\_at} (C++20)}

Od C++20 można użyć funkcji z nagłówka \texttt{<memory>} do
jawnego rozpoczęcia i zakończenia czasu życia obiektu w odpowiednio
przygotowanej pamięci.

\begin{lstlisting}
#include <cstddef>
#include <memory>

struct Tracked {
    int value;
    explicit Tracked(int n) : value(n) {}
};

alignas(Tracked) std::byte buffer[sizeof(Tracked)];
Tracked* obj = std::construct_at(
    reinterpret_cast<Tracked*>(buffer), 42);
// ... korzystamy z *obj ...
std::destroy_at(obj);
\end{lstlisting}

\textbf{Ważne:} \texttt{std::destroy\_at} kończy czas życia obiektu, ale
nie zwalnia jego pamięci. Nie wolno wykonywać dwóch destrukcji tego
samego żywego obiektu ani odwoływać się do obiektu po destrukcji.
Ponadto rozmiar i wyrównanie bufora muszą odpowiadać konstruowanemu
typowi; sam \texttt{reinterpret\_cast} nie naprawia niewłaściwego
wyrównania.

\subsection{Wyrównany \texttt{operator new}}

Dla typów o szczególnie silnym wyrównaniu istnieją przeciążenia
alokacji i dealokacji wykorzystujące \texttt{std::align\_val\_t}
(C++17). Zwykłe \texttt{new Overaligned{}} wybiera odpowiedni
mechanizm automatycznie.

\begin{lstlisting}
#include <new>

struct alignas(128) Block { char data[128]; };
Block* b = new Block{};
delete b;
\end{lstlisting}

Jeśli ręcznie używamy \texttt{::operator new} z argumentem
\texttt{std::align\_val\_t\{128\}}, musimy dobrać zgodną operację
\texttt{::operator delete} z parametrem wyrównania. Mieszanie tych
rodzin funkcji jest błędem.

\subsection{Ponowne użycie tej samej pamięci i \texttt{std::launder}}

W tej samej pamięci można zakończyć życie jednego obiektu, a potem
rozpocząć życie kolejnego, ale dostęp przez wcześniej utworzone
wskaźniki podlega szczegółowym regułom języka. \texttt{std::launder}
(C++17) ma zastosowanie w określonych sytuacjach, gdy nowy obiekt
nie jest przezroczystym zamiennikiem poprzedniego; nie jest
uniwersalnym remedium na aliasing, nieprawidłowy czas życia
czy źle wyrównaną pamięć.

\paragraph{Reguła praktyczna.}
W kodzie aplikacyjnym nie sięgaj po \texttt{std::launder} bez
konkretnego problemu z czasem życia i bez zrozumienia odpowiednich
reguł. Preferuj własność zasobu wyrażoną typem RAII.

\section{Alokatory: pamięć, obiekty i bezpieczeństwo wyjątkowe}
\label{sec:cpp10-allocators}

\subsection{Standardowy model alokatora}

Kontenery standardowe używają \emph{alokatorów} do uzyskiwania
pamięci na swoje elementy. Alokator nie jest tym samym co funkcja
\texttt{new} i nie musi budować obiektu w chwili rezerwacji pamięci.
\texttt{std::allocator\_traits} zapewnia jednolity interfejs do
alokowania, konstruowania i niszczenia obiektów.

\begin{lstlisting}
#include <memory>

struct Item { explicit Item(int v) : value(v) {} int value; };
std::allocator<Item> alloc;
Item* p = alloc.allocate(2); // pamiec na dwa obiekty

std::allocator_traits<decltype(alloc)>::construct(alloc, p, 5);
std::allocator_traits<decltype(alloc)>::construct(alloc, p + 1, 7);

// teraz dwa obiekty sa zywe
std::allocator_traits<decltype(alloc)>::destroy(alloc, p + 1);
std::allocator_traits<decltype(alloc)>::destroy(alloc, p);
alloc.deallocate(p, 2);
\end{lstlisting}

Przykład dla przejrzystości zakłada, że oba konstruktory nie rzucają
wyjątków. \textbf{Pełna implementacja kontenera} musi zliczać
pomyślnie skonstruowane obiekty i w razie wyjątku zniszczyć tylko je,
a potem zwolnić pamięć. Właśnie z tego powodu lepiej opierać się na
gotowych kontenerach niż ręcznie pisać alokację elementów.

\subsection{Rezerwacja pojemności i realokacja \texttt{vector}}

\texttt{std::vector} przechowuje elementy w pamięci ciągłej, dzięki
czemu zapewnia dobrą lokalność dostępu. \texttt{reserve} może
ograniczyć liczbę realokacji, jeśli znamy przewidywaną wielkość
kontenera.

\begin{lstlisting}
#include <vector>

std::vector<int> values;
values.reserve(10'000);
for (int i = 0; i < 10'000; ++i)
    values.push_back(i);
\end{lstlisting}

\texttt{reserve(n)} nie zmienia \texttt{size()}, lecz zapewnia
\texttt{capacity() >= n}, chyba że zgłoszony zostanie wyjątek.
Realokacja unieważnia wskaźniki, referencje i iteratory do elementów.
Nie jest dobrym pomysłem wywoływanie \texttt{reserve(size()+1)}
przy każdym dodaniu elementu: można w ten sposób utracić
korzystną politykę geometrycznego wzrostu pojemności.

\subsection{Dlaczego własny alokator to nie pierwszy krok?}

Własny alokator w szablonie kontenera wiąże się z wymaganiami
dotyczącymi propagowania przy kopiowaniu, przenoszeniu i zamianie
kontenerów oraz z kwestią porównywania alokatorów.
Dla większości zastosowań wygodniejsza jest warstwa
\texttt{std::pmr}, którą poznamy teraz.

\section{\texttt{std::pmr}: polimorficzne zasoby pamięci}

\begin{center}\small
\begin{tabularx}{\textwidth}{@{}p{0.37\textwidth}X@{}}
\toprule
Mechanizm & Zastosowanie i ograniczenie \\
\midrule
\code{pmr::memory_resource} & Interfejs strategii przydziału i zwalniania pamięci używany przez alokatory PMR. \\
\code{pmr::vector<T>} & Kontener, którego alokator korzysta z wybranego zasobu pamięci. \\
\code{monotonic_buffer_resource} & Przydziela kolejne bloki szybko; pojedyncze zwolnienia nie odzyskują miejsca, zasób zwalnia je zbiorczo. \\
\code{unsynchronized_pool_resource} & Pula małych bloków bez synchronizacji między wątkami. \\
\code{synchronized_pool_resource} & Pula z synchronizacją, odpowiednia do współdzielenia między wątkami. \\
\bottomrule
\end{tabularx}
\end{center}
\label{sec:cpp10-pmr}

\subsection{Alokator jako decyzja w czasie działania}

\texttt{std::pmr} (C++17, \texttt{<memory\_resource>}) oddziela
\emph{rodzaj kontenera} od \emph{sposobu przydzielania pamięci}.
\texttt{std::pmr::vector<T>} jest aliasem \texttt{std::vector<T, ...>}
z alokatorem polimorficznym. Zasób \texttt{memory\_resource}
można wybierać w czasie działania programu.

\begin{lstlisting}
#include <memory_resource>
#include <vector>

std::pmr::monotonic_buffer_resource arena;
std::pmr::vector<int> numbers{&arena};
for (int i = 0; i < 100; ++i)
    numbers.push_back(i);
\end{lstlisting}

Kontener przechowuje \emph{wskaźnik do zasobu}, ale nie przejmuje
własności zasobu pamięci. Zasób musi
pozostać żywy wystarczająco długo, aby obsłużyć wszystkie powiązane
alokacje i dealokacje. \textbf{Nie wolno niszczyć lub resetować
areny, gdy żywe obiekty wciąż korzystają z jej pamięci}.

\subsection{Arena monotoniczna: wiele alokacji, jedno zwolnienie}

\texttt{std::pmr::monotonic\_buffer\_resource} wydaje kolejne
fragmenty pamięci. Jej \texttt{deallocate} nie zwalnia pojedynczych
bloków; odzyskanie następuje przy \texttt{release()} lub destrukcji
zasobu. To pasuje do wielu obiektów budowanych razem i usuwanych
w jednej fazie, np. drzewa parsowania pojedynczego dokumentu.

\begin{lstlisting}
#include <cstddef>
#include <memory_resource>
#include <vector>

alignas(std::max_align_t) std::byte storage[4096];
std::pmr::monotonic_buffer_resource arena(
    storage, sizeof(storage),
    std::pmr::null_memory_resource());

{
    std::pmr::vector<int> data{&arena};
    for (int i = 0; i < 32; ++i) data.push_back(i);
} // kontener znika; jego deallocate nie uwalnia blokow
// arena jest teraz bezpieczna do zniszczenia lub release()
\end{lstlisting}

\texttt{null\_memory\_resource()} wymusza niepowodzenie alokacji,
gdy lokalny bufor się wyczerpie (rzucane jest \texttt{std::bad\_alloc}).
Jeżeli zamiast niego użyjemy zwykłego zasobu nadrzędnego
(\emph{upstream}), arena może dobierać kolejne bufory.

\paragraph{Ważna konsekwencja.}
Gdy \texttt{vector} zwiększa pojemność, porzuca stary blok.
W przypadku areny monotonicznej ten blok zwykle pozostaje zajęty
aż do końca całej fazy. Arena nie oznacza zatem automatycznie
najmniejszego zużycia pamięci; warto stosować \texttt{reserve}.

\subsection{Zasoby pulowe: \texttt{unsynchronized\_pool\_resource}}

Gdy obiekty są często alokowane i usuwane w różnych momentach,
przydatne bywają pule przechowujące bloki o podobnych rozmiarach.

\begin{lstlisting}
#include <memory_resource>
#include <vector>

std::pmr::unsynchronized_pool_resource pool;
std::pmr::vector<int> items{&pool};
items.resize(500);
\end{lstlisting}

\texttt{unsynchronized\_pool\_resource} nie zapewnia bezpieczeństwa
jednoczesnego użycia z wielu wątków. Do współdzielonego użycia
istnieje \texttt{synchronized\_pool\_resource}. Wybór zależy od
profilu alokacji i pomiarów: pule nie są uniwersalnie szybsze.

\subsection{Własny zasób: liczymy rzeczywiste alokacje}

Wariant diagnostyczny może zliczać żądania pamięci i delegować je
do \texttt{new\_delete\_resource()}. Trzy wirtualne funkcje
\texttt{do\_allocate}, \texttt{do\_deallocate}, \texttt{do\_is\_equal}
tworzą podstawę własnej implementacji.

\begin{lstlisting}
#include <cstddef>
#include <memory_resource>

struct CountingResource : std::pmr::memory_resource {
    std::pmr::memory_resource* upstream =
        std::pmr::new_delete_resource();
    std::size_t allocations = 0;
    std::size_t bytes = 0;
private:
    void* do_allocate(std::size_t n,
                      std::size_t alignment) override {
        void* p = upstream->allocate(n, alignment);
        ++allocations;
        bytes += n;
        return p;
    }
    void do_deallocate(void* p, std::size_t n,
                       std::size_t alignment) override {
        upstream->deallocate(p, n, alignment);
    }
    bool do_is_equal(
        const std::pmr::memory_resource& other
    ) const noexcept override {
        return this == &other;
    }
};
\end{lstlisting}

\begin{lstlisting}
CountingResource stats;
{
    std::pmr::vector<int> data{&stats};
    for (int i = 0; i < 100; ++i) data.push_back(i);
}
// stats.allocations i stats.bytes opisuja zadania
// przeslane do tego konkretnego zasobu.
\end{lstlisting}

\textbf{Uwaga:} nie są to wszystkie alokacje programu, a
\texttt{bytes} oznacza sumę żądanych rozmiarów, nie największe
jednoczesne zużycie pamięci. Warto liczyć także dealokacje,
\emph{peak usage} i rozkład rozmiarów bloków.

\subsection{Zagnieżdżone kontenery i propagacja zasobu}

\begin{lstlisting}
#include <memory_resource>
#include <string>
#include <vector>

CountingResource stats;
std::pmr::vector<std::pmr::string> names{&stats};
names.emplace_back("alpha");
names.emplace_back("a rather long text that needs allocation");
\end{lstlisting}

\texttt{pmr::string} również korzysta z polimorficznego alokatora,
ale dokładne zachowanie małych napisów zależy od implementacji
(\emph{small string optimization}). Nie należy zakładać, że każdy
\texttt{emplace\_back} przydziela pamięć ani że krótki napis wymaga
alokacji dynamicznej.

\section{Lokalność danych i hierarchia cache}
\label{sec:cpp10-cache}

\subsection{Dlaczego asymptotyka nie wystarcza?}

Dwa algorytmy o złożoności \(O(n)\) mogą mieć bardzo różną wydajność,
jeżeli pierwszy czyta sąsiadujące elementy tablicy, a drugi podąża
za wskaźnikami do rozproszonych obiektów. Procesor pobiera dane
porcjami (\emph{cache lines}); dostęp sekwencyjny sprzyja
prefetchingowi, zaś losowy dostęp bywa ograniczony opóźnieniem pamięci.

Nie ma jednego rozmiaru linii cache dla wszystkich procesorów ani
jednej wartości kosztu trafienia i chybienia. Zależą one od sprzętu,
aktualnego obciążenia i rodzaju dostępu.

\begin{lstlisting}
#include <numeric>
#include <vector>

std::vector<int> v(100000, 1);
long long sum = std::accumulate(v.begin(), v.end(), 0LL);
// Dostep sekwencyjny do pamieci ciaglej.
\end{lstlisting}

Wskaźniki do osobno alokowanych małych węzłów mogą pogorszyć
lokalność i powiększyć narzut pamięciowy. Nie oznacza to,
że \texttt{std::vector} jest zawsze lepszy od \texttt{std::list}:
wybór zależy od wymaganych operacji, unieważniania iteratorów
i rzeczywistego profilu obciążenia.

\subsection{AoS kontra SoA}

W układzie \emph{Array of Structures} (AoS) przechowujemy obok siebie
wszystkie składowe kolejnych rekordów. W \emph{Structure of Arrays}
(SoA) przechowujemy osobne tablice poszczególnych składowych.

\begin{lstlisting}
#include <vector>

struct PointAoS { float x, y, z, intensity; };
std::vector<PointAoS> aos;

struct PointSoA {
    std::vector<float> x, y, z, intensity;
};
PointSoA soa;
\end{lstlisting}

Gdy obliczamy wyłącznie sumę współrzędnych \texttt{x}, układ SoA
umożliwia odczytywanie samych potrzebnych danych, a nie całych
rekordów. Ułatwia to też wektoryzację. Jeśli natomiast w każdej
operacji potrzebujemy wszystkich pól jednego rekordu, AoS może być
korzystniejszy. Wydajność zależy od rozmiaru danych, kodu i sprzętu.

\begin{center}
\small
\begin{tabular}{|p{6.0cm}|p{6.0cm}|}
\hline
\textbf{AoS: rekordy obok siebie} & \textbf{SoA: kolumny obok siebie} \\
\hline
\texttt{[x1 y1 z1 i1] [x2 y2 z2 i2]} &
\texttt{[x1 x2 ...] [y1 y2 ...]} \\
\hline
Wygodny dostęp do całego punktu & Wygodny dostęp do jednej cechy wszystkich punktów \\
\hline
Często mniej kontenerów i prostszy model & Często lepsze wykorzystanie cache w obliczeniach kolumnowych \\
\hline
\end{tabular}
\end{center}

\subsection{Stride, skoki i przewidywanie dostępu}

Przechodzenie po pamięci ze skokiem \(k\) oznacza, że kolejne
operacje mogą dotykać odległych linii cache. Dla macierzy w układzie
wierszowym zwykle korzystniej jest przechodzić po kolumnach
\emph{wewnątrz pojedynczego wiersza}, a następnie po wierszach.
Zanim zmienimy kolejność pętli, należy uwzględnić zależności
obliczeniowe, autowektoryzację i to, czy kompilator już dokonał
danej transformacji.

\begin{lstlisting}
#include <vector>

constexpr std::size_t rows = 100;
constexpr std::size_t cols = 200;
std::vector<double> matrix(rows * cols, 1.0);
double sum = 0.0;
for (std::size_t i = 0; i < rows; ++i)
    for (std::size_t j = 0; j < cols; ++j)
        sum += matrix[i * cols + j];
\end{lstlisting}

\section{False sharing i współbieżne liczniki}
\label{sec:cpp10-false-sharing}

\subsection{Na czym polega false sharing?}

Jeżeli dwa wątki wielokrotnie zapisują \emph{różne} zmienne, ale
zmienne znajdują się w tej samej linii cache, protokół spójności
pamięci podręcznej może wymuszać kosztowną wymianę własności tej
linii między rdzeniami. Nazywamy to \emph{false sharing}.
Nie oznacza to \emph{data race}, jeżeli poszczególne zmienne są
atomowe lub inaczej poprawnie synchronizowane.

\begin{lstlisting}
#include <atomic>
#include <cstdint>

struct NaiveCounters {
    std::atomic<std::uint64_t> first{0};
    std::atomic<std::uint64_t> second{0};
};
\end{lstlisting}

Dwa liczniki mogą być blisko siebie, jednak dokładny layout obiektu
i mapowanie do linii cache nie są gwarantowane. Na niektórych
platformach ich równoczesna modyfikacja będzie kosztowna.

\subsection{Rozdzielenie liczników w pamięci}

\begin{lstlisting}
#include <atomic>
#include <cstdint>

struct alignas(64) IsolatedAtomic {
    std::atomic<std::uint64_t> value{0};
};
struct SeparatedCounters {
    IsolatedAtomic first, second;
};
\end{lstlisting}

Tu użyliśmy 64 bajtów wyłącznie jako \emph{przykładowej} wartości.
Od C++17 w nagłówku \texttt{<new>} istnieją dwie stałe
opisujące zalecane odstępy dla danych często modyfikowanych przez
różne wątki (\emph{destructive interference}) oraz dla danych,
które powinny pozostawać blisko siebie (\emph{constructive
interference}). Pełną nazwę pierwszej z nich pokazuje kod poniżej.
Wartości stałych zależą od implementacji; nie stanowią pomiaru
rozmiaru linii cache w uruchomionym procesorze.

\begin{lstlisting}
#include <atomic>
#include <new>

struct alignas(
    std::hardware_destructive_interference_size
) Counter {
    std::atomic<unsigned> value{0};
};
\end{lstlisting}

Rozdzielenie liczników zwiększa rozmiar struktury. Przy milionie
liczników koszt pamięciowy może przewyższyć korzyść. Technika ma
największy sens, gdy oddzielne wątki często zapisują niewielką
liczbę niezależnych pól.

\subsection{False sharing a semantyka atomowa}

\texttt{std::memory\_order\_relaxed} wystarczy do niektórych
niezależnych liczników, gdy zależy nam wyłącznie na atomowym
zliczaniu; nie gwarantuje jednak synchronizacji innych danych.
Usunięcie false sharing nie zastępuje poprawnego modelu pamięci
współbieżnej. Szczegółowe zasady operacji atomowych były omawiane
w rozdziale~1 o współbieżności.

\section{Jak rzetelnie mierzyć wydajność?}
\label{sec:cpp10-benchmark}

\subsection{Zasady pomiarów}

Najważniejsze pytanie brzmi nie „czy ten kod wygląda szybciej?”,
lecz „czy w rzeczywistym scenariuszu uzyskujemy powtarzalną poprawę?”.
Doświadczenie powinno uwzględniać:
\begin{enumerate}
\item kompilację zoptymalizowaną (np. \texttt{-O2} lub \texttt{-O3});
\item \emph{warm-up} i wielokrotne uruchomienia;
\item jednakowe dane oraz taki sam rezultat dla obu wariantów;
\item ochronę przed usunięciem testowanej pracy przez optymalizator;
\item pomiar rozkładu wyników, a nie pojedynczego odczytu;
\item ograniczenie wpływu alokacji, I/O, schedulerów i skalowania CPU;
\item kontrolę zużycia pamięci, a nie tylko czasu;
\item narzędzia profilujące (np. Linux \texttt{perf}, odpowiedniki
      sprzętowe lub profilery danego środowiska).
\end{enumerate}

\subsection{Minimalna funkcja pomiarowa}

\begin{lstlisting}
#include <chrono>
#include <utility>

template<class F>
auto measure(F&& f) {
    auto start = std::chrono::steady_clock::now();
    auto result = std::forward<F>(f)();
    auto stop = std::chrono::steady_clock::now();
    return std::pair{result, stop - start};
}
\end{lstlisting}

Taka funkcja przydaje się do nauki, ale \emph{nie jest kompletna}
biblioteka benchmarkowa. Dla zbyt krótkich fragmentów narzut
pomiaru może przeważać nad czasem pracy. Dla liczb zmiennoprzecinkowych
porównywanie różnych kolejności sumowania wymaga ostrożności,
ponieważ zmiana kolejności działań może zmienić zaokrąglenia.

\subsection{Sanitizatory i wydajność}

Sanitizatory, np. AddressSanitizer i UndefinedBehaviorSanitizer,
są narzędziami do badania poprawności pamięci i innych błędów.
Wprowadzają narzut, dlatego \emph{pomiarów porównawczych wydajności}
nie należy bezrefleksyjnie wykonywać na buildach sanitizowanych.
Bezpieczny cykl pracy wygląda następująco: test poprawności,
profilowanie zwykłego buildu, zmiana jednej rzeczy, ponowny test
poprawności i ponowny pomiar.

\section{Przykładowy projekt: szybkie przetwarzanie wielu punktów}
\label{sec:cpp10-project}

Przyjmijmy, że program przechowuje milion punktów \((x,y,z,I)\)
i dla każdego oblicza tylko prosty filtr zależny od \(x\) i \(I\).
Możemy rozważyć dwa modele:
\begin{enumerate}
\item \textbf{AoS:} \texttt{vector<Point>} --- prostszy interfejs,
      wygodny dostęp do pełnego punktu.
\item \textbf{SoA:} osobne \texttt{vector<float>} na \(x,y,z,I\) ---
      możliwy mniejszy transfer danych potrzebnych do filtru.
\end{enumerate}

Następnie oceniamy, czy warto użyć \texttt{reserve} przed zapełnieniem
kontenerów; jeśli rekordy są tworzone na potrzeby jednego zapytania,
możemy sprawdzić arenę monotoniczną. Trzeba jednak rozróżnić
\emph{koszt wypełnienia struktur} od \emph{kosztu przetwarzania}: arena
zmienia zachowanie alokacji, nie gwarantuje poprawy pracy procesora
nad danymi.

Dla kodu numerycznego duży wpływ może mieć też wybór kolejności
wymiarów macierzy, reprezentacji sąsiedztwa w siatce czy sposobu
obliczania operatora na konkretnym punkcie zamiast konstruowania
dużych obiektów pośrednich. Właściwa reprezentacja danych potrafi
dać większy zysk niż drobne optymalizacje pojedynczych instrukcji.

\section{Jak dobrać technikę?}

\begin{center}
\small
\begin{tabular}{|p{5.7cm}|p{6.7cm}|}
\hline
\textbf{Problem} & \textbf{Pierwszy sensowny krok} \\
\hline
Za dużo dynamicznych alokacji & \texttt{reserve}, ograniczenie obiektów pośrednich, pomiar alokacji \\
\hline
Wiele obiektów znika razem & \texttt{pmr::monotonic\_buffer\_resource} po kontroli czasu życia \\
\hline
Dużo małych alokacji o różnej żywotności & \texttt{pmr::...pool\_resource} i pomiar \\
\hline
Praca nad jedną kolumną wielu rekordów & Porównanie AoS z SoA \\
\hline
Losowe skoki po pamięci & Sprawdzenie lokalności, zmian formatu danych i blokowania pętli \\
\hline
Niezależne atomiki spowalniają się nawzajem & Analiza false sharing i separacja intensywnie zapisywanych pól \\
\hline
Trzeba uzyskać silne wyrównanie & \texttt{alignas}, odpowiednia alokacja i kontrola \texttt{alignof} \\
\hline
Trzeba napisać kontener lub pulę & \texttt{allocator\_traits}, zarządzanie czasem życia, RAII \\
\hline
Potrzeba przenośnej serializacji & Jawny format bajtowy, bez kopiowania całych struktur \\
\hline
\end{tabular}
\end{center}

\section{Typowe błędy i pułapki}

\begin{enumerate}
\item \textbf{Wniosek „\texttt{sizeof(T)} to suma pól”.}
      Padding może zwiększyć rozmiar, a reguły klas mogą być złożone.
\item \textbf{Uznanie layoutu z jednej kompilacji za standard.}
      ABI, wyrównanie i rozmiary typów różnią się między platformami.
\item \textbf{Nieprawidłowe \texttt{reinterpret\_cast}.}
      Sama zmiana typu wskaźnika nie legalizuje dostępu do danych.
\item \textbf{Mylenie pamięci z czasem życia obiektu.}
      \texttt{allocate} nie zastępuje konstruowania obiektów klasowych.
\item \textbf{Zwalnianie pamięci bez destruktora.}
      Dla nietrywialnych typów może wyciekać zasób zarządzany przez obiekt.
\item \textbf{Niezgodna para alokacja--dealokacja.}
      Szczególnie ważne przy over-alignment i ręcznej alokacji.
\item \textbf{Pozostawienie kontenera po zniszczeniu areny \texttt{pmr}.}
      Jego pamięć staje się niepoprawna; wynik to niezdefiniowane zachowanie.
\item \textbf{Wywołanie \texttt{release()} przy żywych obiektach.}
      Arena może unieważnić ich pamięć bez wywoływania destruktorów.
\item \textbf{Założenie, że \texttt{monotonic\_buffer\_resource} oszczędza
      pamięć.} Zwolnienie pojedynczych bloków zwykle nie następuje.
\item \textbf{Współdzielenie \texttt{unsynchronized\_pool\_resource}
      między wątkami bez ochrony.} Zasób nie jest thread-safe.
\item \textbf{Optymalizacja poprzez mniejszą strukturę bez pomiaru.}
      Mniejszy \texttt{sizeof} nie gwarantuje krótszego czasu wykonania.
\item \textbf{Rzekoma optymalizacja przez \texttt{volatile}.}
      \texttt{volatile} nie zapewnia synchronizacji między wątkami.
\item \textbf{Pomiar jednego przebiegu.} Scheduler, częstotliwości
      zegara i cache mogą silnie zniekształcić wynik.
\item \textbf{Użycie sanitizatora jako miary szybkości.}
      Narzut instrumentacji zmienia charakterystykę programu.
\end{enumerate}

\paragraph{Pytania kontrolne.}
\begin{enumerate}
\item Czym różnią się \texttt{sizeof(T)} i \texttt{alignof(T)}?
\item Dlaczego zmiana kolejności pól może zmienić rozmiar struktury?
\item Czy placement \texttt{new} przydziela pamięć?
\item Jaka jest różnica między \texttt{destroy\_at} a \texttt{deallocate}?
\item Kiedy arena monotoniczna może zużywać więcej pamięci niż oczekujemy?
\item Dlaczego \texttt{pmr::vector<T>} nie jest właścicielem zasobu?
\item Kiedy SoA może przyspieszyć obliczenia na rekordach?
\item Czy false sharing i data race są tym samym problemem?
\item Dlaczego \texttt{std::memory\_order\_relaxed} nie rozwiązuje
      wszystkich problemów synchronizacji?
\item Jak odróżnić realną poprawę wydajności od szumu pomiarowego?
\end{enumerate}

\paragraph{Dalsza lektura.}
Dokumentacja \emph{cppreference} obejmuje hasła \emph{Object},
\emph{Object lifetime}, \texttt{alignof}, \texttt{alignas},
\texttt{std::construct\_at} oraz \texttt{allocator\_traits}.
Osobno warto przeczytać opis zasobów \texttt{std::pmr}
i własności \emph{hardware interference size}.
W praktyce polecam analizę śladów alokacji, profilery
wydajności i \emph{Compiler Explorer} do porównywania
wygenerowanego kodu maszynowego.


\chapter[Organizacja pamięci i wydajność]{Organizacja pamięci, wyrównanie i wydajność niskopoziomowa w C++}
\label{chap:cpp-memory-performance}

\section{Po co analizować organizację pamięci?}

W rozdziałach o wskaźnikach (rozdział~3), semantyce przenoszenia
(rozdział~6) i iteratorach (rozdział~8) interesowało nas głównie
\emph{poprawne zarządzanie obiektami}. Teraz przyjrzymy się temu, jak
obiekty są reprezentowane w pamięci oraz w jaki sposób wybory
programisty wpływają na liczbę alokacji i lokalność dostępu do danych.

Rozróżniajmy trzy poziomy opisu:
\begin{enumerate}
\item \textbf{Gwarancje języka C++:} rozmiar obiektu, wymagane wyrównanie,
      czas życia obiektu, legalność dostępu i zasady alokowania pamięci.
\item \textbf{Szczegóły implementacji:} rozmieszczenie pól, rozmiary
      wskaźników, zachowanie konkretnego alokatora czy układ binarny ABI.
\item \textbf{Właściwości sprzętu:} linie cache, hierarchia pamięci,
      prefetching, przepustowość i zjawisko \emph{false sharing}.
\end{enumerate}

Kod poprawny według standardu może być wolny; szybki kod działający
na jednym komputerze może być powolny na innym. Dlatego optymalizacja
wymaga najpierw \emph{pomiaru}, a następnie weryfikacji hipotezy.

\section{Obiekt, bajt, adres i czas życia}
\label{sec:cpp10-object}

\subsection{Pamięć nie jest tym samym co obiekt}

W modelu C++ \emph{obszar pamięci} może istnieć bez żywego obiektu typu
\texttt{T}. Obiekt charakteryzują między innymi typ, rozmiar, czas życia
i, dla pewnych klas, odpowiednio wywołane konstruktory i destruktory.
Samo uzyskanie adresu wystarczająco dużego bloku nie oznacza, że
wolno traktować ten blok jak dowolny obiekt C++.

\begin{lstlisting}
#include <cstddef>
#include <type_traits>

static_assert(sizeof(char) == 1);
static_assert(sizeof(std::byte) == 1);
static_assert(sizeof(int) >= 1);

struct Point { double x, y; };
static_assert(sizeof(Point) >= 2 * sizeof(double));
\end{lstlisting}

\texttt{sizeof(T)} podaje rozmiar obiektu w \emph{bajtach C++}.
Jeden bajt nie musi mieć ośmiu bitów; ich liczbę określa
\texttt{CHAR\_BIT} z nagłówka \texttt{<climits>}.
Rozmiar typu obiektowego jest dodatni. Każdy obiekt ma też
wymagane \emph{wyrównanie}, które omówimy dalej.

\paragraph{Czas przechowywania a czas życia.}
Należy odróżnić \emph{storage duration} od \emph{object lifetime}:
\begin{itemize}
\item obiekt lokalny ma zazwyczaj automatyczny czas przechowywania;
\item obiekt z \texttt{static} ma statyczny czas przechowywania;
\item obiekt \texttt{thread\_local} ma czas przechowywania wątku;
\item pamięć otrzymana przez \texttt{new} lub alokator jest pamięcią
      o dynamicznym czasie przechowywania.
\end{itemize}
Nie należy utożsamiać tych kategorii z gwarantowanymi fizycznymi
obszarami „stos” i „sterta”: te nazwy opisują zwyczajowe implementacje,
a nie pełny model abstrakcyjny języka.

\begin{lstlisting}
#include <memory>

int global = 17;                 // static storage duration
thread_local int per_thread = 0; // thread storage duration

void demo() {
    int local = 5;               // automatic storage duration
    auto p = std::make_unique<int>(10);
    // Obiekt *p zyje w pamieci zarzadzanej dynamicznie.
}
\end{lstlisting}

\subsection{Reprezentacja obiektu i bezpieczne oglądanie bajtów}

C++ pozwala badać reprezentację bajtową obiektu przez
\texttt{char}, \texttt{unsigned char} lub \texttt{std::byte}.
Nie wolno jednak wyciągać z tego wniosku, że dowolny wskaźnik
\texttt{T*} można przekonwertować na \texttt{U*} i odczytać dane jako
\texttt{U}: takie postępowanie może naruszyć reguły wyrównania,
czasu życia i dostępu przez glvalue odpowiedniego typu.

\begin{lstlisting}
#include <array>
#include <cstddef>
#include <cstring>
#include <cstdint>

std::uint32_t x = 0x12345678;
std::array<std::byte, sizeof(x)> bytes{};
std::memcpy(bytes.data(), &x, sizeof(x));

std::uint32_t y = 0;
std::memcpy(&y, bytes.data(), sizeof(y));
// y == x
\end{lstlisting}

Kolejność bajtów zależy od \emph{endianness}. Można ją rozpoznać przez
\texttt{std::endian} (C++20, nagłówek \texttt{<bit>}), jednak do
serializacji najlepiej jawnie ustalać porządek bajtów, format pól
i wersję protokołu. Nie należy wysyłać całej struktury przez
\texttt{write(fd, \&object, sizeof(object))} jako przenośnego formatu:
padding, reprezentacja liczb i ABI mogą się różnić.

\subsection{\texttt{std::bit\_cast} i ograniczenia rzutowań}

\texttt{std::bit\_cast<To>(from)} (C++20, \texttt{<bit>}) tworzy wartość
\texttt{To} przez skopiowanie reprezentacji bitowej, o ile oba typy
są trywialnie kopiowalne i mają równy rozmiar. Dodatkowo reprezentacja
musi oznaczać poprawną wartość docelowego typu.

\begin{lstlisting}
#include <bit>
#include <cstdint>

float f = 1.0f;
// Przyklad zaklada sizeof(float) == sizeof(uint32_t).
// Na typowej platformie wynik to kod bitowy IEEE 754.
if constexpr (sizeof(float) == sizeof(std::uint32_t)) {
    auto bits = std::bit_cast<std::uint32_t>(f);
    (void)bits;
}
\end{lstlisting}

Nie mylmy \texttt{bit\_cast} z \texttt{reinterpret\_cast}:
pierwszy wytwarza \emph{wartość}, drugi może tylko zmienić sposób
interpretacji wskaźnika; sam nie tworzy obiektu docelowego.

\section{Layout struktury: \texttt{sizeof}, \texttt{alignof}, padding}

\begin{center}\small
\begin{tabularx}{\textwidth}{@{}p{0.33\textwidth}X@{}}
\toprule
Narzędzie & Co pokazuje lub zmienia \\
\midrule
\code{sizeof(T)} & Podaje rozmiar obiektu typu \code{T} w bajtach, wraz z ewentualnym paddingiem. \\
\code{alignof(T)} & Podaje wymagane wyrównanie adresu obiektu typu \code{T}. \\
\code{offsetof(T, field)} & Podaje przesunięcie składowej w typie o standardowym układzie. \\
\code{alignas(n)} & Żąda wyrównania co najmniej do granicy \code{n} dla obiektu lub typu. \\
\code{std::byte} & Reprezentuje jednostkę surowej pamięci bez znaczenia znaku tekstowego. \\
\bottomrule
\end{tabularx}
\end{center}
\label{sec:cpp10-layout}

\subsection{Skąd bierze się padding?}

Każdy typ obiektowy ma wymaganie wyrównania. Kompilator może wstawić
\emph{padding}, czyli nieużywane bajty, aby kolejne pola miały
odpowiednie adresy, a elementy tablicy struktur były poprawnie
wyrównane. Rozmieszczenie pól zależy od typu klasy i reguł języka;
konkretna liczba bajtów paddingu zależy od implementacji.

\begin{lstlisting}
#include <cstddef>
#include <iostream>

struct Unordered {
    char flag;
    double value;
    int count;
};

struct Ordered {
    double value;
    int count;
    char flag;
};

std::cout << sizeof(Unordered) << '\n';
std::cout << sizeof(Ordered) << '\n';
std::cout << alignof(Ordered) << '\n';
\end{lstlisting}

Na częstym ABI, przy \texttt{sizeof(double)==8},
\texttt{alignof(double)==8} i \texttt{sizeof(int)==4},
pierwszy typ zajmuje zwykle 24 bajty, a drugi 16. \textbf{Nie są to
wartości wymagane przez standard}. Samo przestawienie pól może
zmniejszyć zapotrzebowanie na pamięć milionów rekordów, ale może też
pogorszyć czytelność, stabilność ABI lub lokalność dostępu.

\begin{center}
\small
\begin{tabular}{|l|p{4.4cm}|p{4.4cm}|}
\hline
 & \textbf{Unordered (przykładowo)} & \textbf{Ordered (przykładowo)} \\
\hline
Początek & \texttt{flag}: 1 B & \texttt{value}: 8 B \\
\hline
Dalej & padding: 7 B & \texttt{count}: 4 B \\
\hline
Dalej & \texttt{value}: 8 B & \texttt{flag}: 1 B \\
\hline
Koniec & \texttt{count}: 4 B + padding: 4 B & padding: 3 B \\
\hline
Razem & 24 B & 16 B \\
\hline
\end{tabular}
\end{center}

\subsection{\texttt{offsetof}, standard-layout i ABI}

Dla typów \emph{standard-layout} makro \texttt{offsetof} pozwala
poznać przesunięcie pola od początku obiektu. Używanie go z innymi
typami klas od C++17 jest warunkowo wspierane, dlatego do analiz
przenośnych przyjmujemy tutaj standard-layout.

\begin{lstlisting}
#include <cstddef>
#include <cstdint>
#include <type_traits>

struct Header {
    std::uint32_t id;
    std::uint16_t version;
};
static_assert(std::is_standard_layout_v<Header>);
static_assert(offsetof(Header, version) >=
              sizeof(std::uint32_t));
\end{lstlisting}

W analizie ABI trzeba pamiętać o polach wirtualnych, dziedziczeniu,
przestrzeniach adresowych i konwencjach wywołania. Zmiana kolejności
pól publicznej struktury używanej w bibliotekach binarnych może
zmienić jej ABI nawet wtedy, gdy program źródłowy nadal się kompiluje.

\subsection{\texttt{alignof} i \texttt{alignas}}

\texttt{alignof(T)} zwraca wyrównanie typu \texttt{T};
\texttt{alignas(N)} pozwala narzucić silniejsze wyrównanie.
Należy zachować ograniczenia języka dotyczące dopuszczalnych wartości
\texttt{N}. Wyrównanie może zwiększyć rozmiar obiektu, co jest
istotne przy tablicach.

\begin{lstlisting}
#include <atomic>
#include <cstdint>

struct alignas(64) AlignedCounter {
    std::atomic<std::uint64_t> value{0};
};
static_assert(alignof(AlignedCounter) >= 64);
\end{lstlisting}

Wyrównanie do 64 bajtów ma sens tylko wtedy, gdy wynika z wymagań
sprzętu lub pomiarów. Nie jest to uniwersalna wielkość linii cache.
Ponadto \texttt{alignas} nie czyni dostępu do obiektu atomowym,
nie zastępuje \texttt{std::atomic} i samo nie zapewnia synchronizacji.

\subsection{Czy warto „pakować” struktury?}

Atrybuty kompilatora w rodzaju \texttt{\_\_attribute\_\_((packed))}
lub \texttt{\#pragma pack} nie stanowią przenośnych mechanizmów
standardowego C++20. Zmniejszenie rozmiaru może wymusić niewyrównany
dostęp, obniżyć wydajność, a na niektórych architekturach być
problemem poprawności. W kodzie protokołów lepiej często dekodować
pola jawnie z bufora bajtów i nie wiązać formatu z layoutem klasy.

\section{Wyrównane alokacje i czas życia obiektów}
\label{sec:cpp10-lifetime}

\subsection{Różnica między alokacją a konstruowaniem}

Klasyczne \texttt{new T(args)} wykonuje logicznie dwa zadania:
pozyskuje odpowiednio wyrównaną pamięć i inicjalizuje w niej obiekt.
Analogicznie \texttt{delete} niszczy obiekt i zwalnia pamięć.
Dla narzędzi niskopoziomowych warto rozdzielić te operacje.

\begin{lstlisting}
#include <new>

struct Widget { explicit Widget(int) {} };

void* mem = ::operator new(sizeof(Widget));
Widget* w = nullptr;
try {
    w = ::new (mem) Widget(5); // placement new
} catch (...) {
    ::operator delete(mem);
    throw;
}
w->~Widget();                 // koniec zycia obiektu
::operator delete(mem);      // zwolnienie samej pamieci
\end{lstlisting}

Taka konstrukcja jest poprawna dla \texttt{Widget} o wyrównaniu
obsługiwanym przez zwykłe \texttt{operator new}. Dla typów
\emph{over-aligned} należy korzystać z odpowiedniej alokacji
obsługującej \texttt{std::align\_val\_t} lub ze standardowego
alokatora, który uwzględnia wymagania wyrównania.

\paragraph{Dlaczego to nie jest zalecany idiom aplikacyjny?}
Obsługa wyjątków, wielokrotnych destruktorów, rozmiaru tablicy
i zgodności alokacji z dealokacją łatwo prowadzi do błędów.
W większości kodu wystarczają \texttt{std::vector},
\texttt{std::make\_unique} i obiekty automatyczne. Rozdzielanie
alokacji od konstruowania ma sens głównie przy implementowaniu
kontenerów, pul pamięci i innych abstrakcji zasobów.

\subsection{\texttt{std::construct\_at} i \texttt{std::destroy\_at} (C++20)}

Od C++20 można użyć funkcji z nagłówka \texttt{<memory>} do
jawnego rozpoczęcia i zakończenia czasu życia obiektu w odpowiednio
przygotowanej pamięci.

\begin{lstlisting}
#include <cstddef>
#include <memory>

struct Tracked {
    int value;
    explicit Tracked(int n) : value(n) {}
};

alignas(Tracked) std::byte buffer[sizeof(Tracked)];
Tracked* obj = std::construct_at(
    reinterpret_cast<Tracked*>(buffer), 42);
// ... korzystamy z *obj ...
std::destroy_at(obj);
\end{lstlisting}

\textbf{Ważne:} \texttt{std::destroy\_at} kończy czas życia obiektu, ale
nie zwalnia jego pamięci. Nie wolno wykonywać dwóch destrukcji tego
samego żywego obiektu ani odwoływać się do obiektu po destrukcji.
Ponadto rozmiar i wyrównanie bufora muszą odpowiadać konstruowanemu
typowi; sam \texttt{reinterpret\_cast} nie naprawia niewłaściwego
wyrównania.

\subsection{Wyrównany \texttt{operator new}}

Dla typów o szczególnie silnym wyrównaniu istnieją przeciążenia
alokacji i dealokacji wykorzystujące \texttt{std::align\_val\_t}
(C++17). Zwykłe \texttt{new Overaligned{}} wybiera odpowiedni
mechanizm automatycznie.

\begin{lstlisting}
#include <new>

struct alignas(128) Block { char data[128]; };
Block* b = new Block{};
delete b;
\end{lstlisting}

Jeśli ręcznie używamy \texttt{::operator new} z argumentem
\texttt{std::align\_val\_t\{128\}}, musimy dobrać zgodną operację
\texttt{::operator delete} z parametrem wyrównania. Mieszanie tych
rodzin funkcji jest błędem.

\subsection{Ponowne użycie tej samej pamięci i \texttt{std::launder}}

W tej samej pamięci można zakończyć życie jednego obiektu, a potem
rozpocząć życie kolejnego, ale dostęp przez wcześniej utworzone
wskaźniki podlega szczegółowym regułom języka. \texttt{std::launder}
(C++17) ma zastosowanie w określonych sytuacjach, gdy nowy obiekt
nie jest przezroczystym zamiennikiem poprzedniego; nie jest
uniwersalnym remedium na aliasing, nieprawidłowy czas życia
czy źle wyrównaną pamięć.

\paragraph{Reguła praktyczna.}
W kodzie aplikacyjnym nie sięgaj po \texttt{std::launder} bez
konkretnego problemu z czasem życia i bez zrozumienia odpowiednich
reguł. Preferuj własność zasobu wyrażoną typem RAII.

\section{Alokatory: pamięć, obiekty i bezpieczeństwo wyjątkowe}
\label{sec:cpp10-allocators}

\subsection{Standardowy model alokatora}

Kontenery standardowe używają \emph{alokatorów} do uzyskiwania
pamięci na swoje elementy. Alokator nie jest tym samym co funkcja
\texttt{new} i nie musi budować obiektu w chwili rezerwacji pamięci.
\texttt{std::allocator\_traits} zapewnia jednolity interfejs do
alokowania, konstruowania i niszczenia obiektów.

\begin{lstlisting}
#include <memory>

struct Item { explicit Item(int v) : value(v) {} int value; };
std::allocator<Item> alloc;
Item* p = alloc.allocate(2); // pamiec na dwa obiekty

std::allocator_traits<decltype(alloc)>::construct(alloc, p, 5);
std::allocator_traits<decltype(alloc)>::construct(alloc, p + 1, 7);

// teraz dwa obiekty sa zywe
std::allocator_traits<decltype(alloc)>::destroy(alloc, p + 1);
std::allocator_traits<decltype(alloc)>::destroy(alloc, p);
alloc.deallocate(p, 2);
\end{lstlisting}

Przykład dla przejrzystości zakłada, że oba konstruktory nie rzucają
wyjątków. \textbf{Pełna implementacja kontenera} musi zliczać
pomyślnie skonstruowane obiekty i w razie wyjątku zniszczyć tylko je,
a potem zwolnić pamięć. Właśnie z tego powodu lepiej opierać się na
gotowych kontenerach niż ręcznie pisać alokację elementów.

\subsection{Rezerwacja pojemności i realokacja \texttt{vector}}

\texttt{std::vector} przechowuje elementy w pamięci ciągłej, dzięki
czemu zapewnia dobrą lokalność dostępu. \texttt{reserve} może
ograniczyć liczbę realokacji, jeśli znamy przewidywaną wielkość
kontenera.

\begin{lstlisting}
#include <vector>

std::vector<int> values;
values.reserve(10'000);
for (int i = 0; i < 10'000; ++i)
    values.push_back(i);
\end{lstlisting}

\texttt{reserve(n)} nie zmienia \texttt{size()}, lecz zapewnia
\texttt{capacity() >= n}, chyba że zgłoszony zostanie wyjątek.
Realokacja unieważnia wskaźniki, referencje i iteratory do elementów.
Nie jest dobrym pomysłem wywoływanie \texttt{reserve(size()+1)}
przy każdym dodaniu elementu: można w ten sposób utracić
korzystną politykę geometrycznego wzrostu pojemności.

\subsection{Dlaczego własny alokator to nie pierwszy krok?}

Własny alokator w szablonie kontenera wiąże się z wymaganiami
dotyczącymi propagowania przy kopiowaniu, przenoszeniu i zamianie
kontenerów oraz z kwestią porównywania alokatorów.
Dla większości zastosowań wygodniejsza jest warstwa
\texttt{std::pmr}, którą poznamy teraz.

\section{\texttt{std::pmr}: polimorficzne zasoby pamięci}

\begin{center}\small
\begin{tabularx}{\textwidth}{@{}p{0.37\textwidth}X@{}}
\toprule
Mechanizm & Zastosowanie i ograniczenie \\
\midrule
\code{pmr::memory_resource} & Interfejs strategii przydziału i zwalniania pamięci używany przez alokatory PMR. \\
\code{pmr::vector<T>} & Kontener, którego alokator korzysta z wybranego zasobu pamięci. \\
\code{monotonic_buffer_resource} & Przydziela kolejne bloki szybko; pojedyncze zwolnienia nie odzyskują miejsca, zasób zwalnia je zbiorczo. \\
\code{unsynchronized_pool_resource} & Pula małych bloków bez synchronizacji między wątkami. \\
\code{synchronized_pool_resource} & Pula z synchronizacją, odpowiednia do współdzielenia między wątkami. \\
\bottomrule
\end{tabularx}
\end{center}
\label{sec:cpp10-pmr}

\subsection{Alokator jako decyzja w czasie działania}

\texttt{std::pmr} (C++17, \texttt{<memory\_resource>}) oddziela
\emph{rodzaj kontenera} od \emph{sposobu przydzielania pamięci}.
\texttt{std::pmr::vector<T>} jest aliasem \texttt{std::vector<T, ...>}
z alokatorem polimorficznym. Zasób \texttt{memory\_resource}
można wybierać w czasie działania programu.

\begin{lstlisting}
#include <memory_resource>
#include <vector>

std::pmr::monotonic_buffer_resource arena;
std::pmr::vector<int> numbers{&arena};
for (int i = 0; i < 100; ++i)
    numbers.push_back(i);
\end{lstlisting}

Kontener przechowuje \emph{wskaźnik do zasobu}, ale nie przejmuje
własności zasobu pamięci. Zasób musi
pozostać żywy wystarczająco długo, aby obsłużyć wszystkie powiązane
alokacje i dealokacje. \textbf{Nie wolno niszczyć lub resetować
areny, gdy żywe obiekty wciąż korzystają z jej pamięci}.

\subsection{Arena monotoniczna: wiele alokacji, jedno zwolnienie}

\texttt{std::pmr::monotonic\_buffer\_resource} wydaje kolejne
fragmenty pamięci. Jej \texttt{deallocate} nie zwalnia pojedynczych
bloków; odzyskanie następuje przy \texttt{release()} lub destrukcji
zasobu. To pasuje do wielu obiektów budowanych razem i usuwanych
w jednej fazie, np. drzewa parsowania pojedynczego dokumentu.

\begin{lstlisting}
#include <cstddef>
#include <memory_resource>
#include <vector>

alignas(std::max_align_t) std::byte storage[4096];
std::pmr::monotonic_buffer_resource arena(
    storage, sizeof(storage),
    std::pmr::null_memory_resource());

{
    std::pmr::vector<int> data{&arena};
    for (int i = 0; i < 32; ++i) data.push_back(i);
} // kontener znika; jego deallocate nie uwalnia blokow
// arena jest teraz bezpieczna do zniszczenia lub release()
\end{lstlisting}

\texttt{null\_memory\_resource()} wymusza niepowodzenie alokacji,
gdy lokalny bufor się wyczerpie (rzucane jest \texttt{std::bad\_alloc}).
Jeżeli zamiast niego użyjemy zwykłego zasobu nadrzędnego
(\emph{upstream}), arena może dobierać kolejne bufory.

\paragraph{Ważna konsekwencja.}
Gdy \texttt{vector} zwiększa pojemność, porzuca stary blok.
W przypadku areny monotonicznej ten blok zwykle pozostaje zajęty
aż do końca całej fazy. Arena nie oznacza zatem automatycznie
najmniejszego zużycia pamięci; warto stosować \texttt{reserve}.

\subsection{Zasoby pulowe: \texttt{unsynchronized\_pool\_resource}}

Gdy obiekty są często alokowane i usuwane w różnych momentach,
przydatne bywają pule przechowujące bloki o podobnych rozmiarach.

\begin{lstlisting}
#include <memory_resource>
#include <vector>

std::pmr::unsynchronized_pool_resource pool;
std::pmr::vector<int> items{&pool};
items.resize(500);
\end{lstlisting}

\texttt{unsynchronized\_pool\_resource} nie zapewnia bezpieczeństwa
jednoczesnego użycia z wielu wątków. Do współdzielonego użycia
istnieje \texttt{synchronized\_pool\_resource}. Wybór zależy od
profilu alokacji i pomiarów: pule nie są uniwersalnie szybsze.

\subsection{Własny zasób: liczymy rzeczywiste alokacje}

Wariant diagnostyczny może zliczać żądania pamięci i delegować je
do \texttt{new\_delete\_resource()}. Trzy wirtualne funkcje
\texttt{do\_allocate}, \texttt{do\_deallocate}, \texttt{do\_is\_equal}
tworzą podstawę własnej implementacji.

\begin{lstlisting}
#include <cstddef>
#include <memory_resource>

struct CountingResource : std::pmr::memory_resource {
    std::pmr::memory_resource* upstream =
        std::pmr::new_delete_resource();
    std::size_t allocations = 0;
    std::size_t bytes = 0;
private:
    void* do_allocate(std::size_t n,
                      std::size_t alignment) override {
        void* p = upstream->allocate(n, alignment);
        ++allocations;
        bytes += n;
        return p;
    }
    void do_deallocate(void* p, std::size_t n,
                       std::size_t alignment) override {
        upstream->deallocate(p, n, alignment);
    }
    bool do_is_equal(
        const std::pmr::memory_resource& other
    ) const noexcept override {
        return this == &other;
    }
};
\end{lstlisting}

\begin{lstlisting}
CountingResource stats;
{
    std::pmr::vector<int> data{&stats};
    for (int i = 0; i < 100; ++i) data.push_back(i);
}
// stats.allocations i stats.bytes opisuja zadania
// przeslane do tego konkretnego zasobu.
\end{lstlisting}

\textbf{Uwaga:} nie są to wszystkie alokacje programu, a
\texttt{bytes} oznacza sumę żądanych rozmiarów, nie największe
jednoczesne zużycie pamięci. Warto liczyć także dealokacje,
\emph{peak usage} i rozkład rozmiarów bloków.

\subsection{Zagnieżdżone kontenery i propagacja zasobu}

\begin{lstlisting}
#include <memory_resource>
#include <string>
#include <vector>

CountingResource stats;
std::pmr::vector<std::pmr::string> names{&stats};
names.emplace_back("alpha");
names.emplace_back("a rather long text that needs allocation");
\end{lstlisting}

\texttt{pmr::string} również korzysta z polimorficznego alokatora,
ale dokładne zachowanie małych napisów zależy od implementacji
(\emph{small string optimization}). Nie należy zakładać, że każdy
\texttt{emplace\_back} przydziela pamięć ani że krótki napis wymaga
alokacji dynamicznej.

\section{Lokalność danych i hierarchia cache}
\label{sec:cpp10-cache}

\subsection{Dlaczego asymptotyka nie wystarcza?}

Dwa algorytmy o złożoności \(O(n)\) mogą mieć bardzo różną wydajność,
jeżeli pierwszy czyta sąsiadujące elementy tablicy, a drugi podąża
za wskaźnikami do rozproszonych obiektów. Procesor pobiera dane
porcjami (\emph{cache lines}); dostęp sekwencyjny sprzyja
prefetchingowi, zaś losowy dostęp bywa ograniczony opóźnieniem pamięci.

Nie ma jednego rozmiaru linii cache dla wszystkich procesorów ani
jednej wartości kosztu trafienia i chybienia. Zależą one od sprzętu,
aktualnego obciążenia i rodzaju dostępu.

\begin{lstlisting}
#include <numeric>
#include <vector>

std::vector<int> v(100000, 1);
long long sum = std::accumulate(v.begin(), v.end(), 0LL);
// Dostep sekwencyjny do pamieci ciaglej.
\end{lstlisting}

Wskaźniki do osobno alokowanych małych węzłów mogą pogorszyć
lokalność i powiększyć narzut pamięciowy. Nie oznacza to,
że \texttt{std::vector} jest zawsze lepszy od \texttt{std::list}:
wybór zależy od wymaganych operacji, unieważniania iteratorów
i rzeczywistego profilu obciążenia.

\subsection{AoS kontra SoA}

W układzie \emph{Array of Structures} (AoS) przechowujemy obok siebie
wszystkie składowe kolejnych rekordów. W \emph{Structure of Arrays}
(SoA) przechowujemy osobne tablice poszczególnych składowych.

\begin{lstlisting}
#include <vector>

struct PointAoS { float x, y, z, intensity; };
std::vector<PointAoS> aos;

struct PointSoA {
    std::vector<float> x, y, z, intensity;
};
PointSoA soa;
\end{lstlisting}

Gdy obliczamy wyłącznie sumę współrzędnych \texttt{x}, układ SoA
umożliwia odczytywanie samych potrzebnych danych, a nie całych
rekordów. Ułatwia to też wektoryzację. Jeśli natomiast w każdej
operacji potrzebujemy wszystkich pól jednego rekordu, AoS może być
korzystniejszy. Wydajność zależy od rozmiaru danych, kodu i sprzętu.

\begin{center}
\small
\begin{tabular}{|p{6.0cm}|p{6.0cm}|}
\hline
\textbf{AoS: rekordy obok siebie} & \textbf{SoA: kolumny obok siebie} \\
\hline
\texttt{[x1 y1 z1 i1] [x2 y2 z2 i2]} &
\texttt{[x1 x2 ...] [y1 y2 ...]} \\
\hline
Wygodny dostęp do całego punktu & Wygodny dostęp do jednej cechy wszystkich punktów \\
\hline
Często mniej kontenerów i prostszy model & Często lepsze wykorzystanie cache w obliczeniach kolumnowych \\
\hline
\end{tabular}
\end{center}

\subsection{Stride, skoki i przewidywanie dostępu}

Przechodzenie po pamięci ze skokiem \(k\) oznacza, że kolejne
operacje mogą dotykać odległych linii cache. Dla macierzy w układzie
wierszowym zwykle korzystniej jest przechodzić po kolumnach
\emph{wewnątrz pojedynczego wiersza}, a następnie po wierszach.
Zanim zmienimy kolejność pętli, należy uwzględnić zależności
obliczeniowe, autowektoryzację i to, czy kompilator już dokonał
danej transformacji.

\begin{lstlisting}
#include <vector>

constexpr std::size_t rows = 100;
constexpr std::size_t cols = 200;
std::vector<double> matrix(rows * cols, 1.0);
double sum = 0.0;
for (std::size_t i = 0; i < rows; ++i)
    for (std::size_t j = 0; j < cols; ++j)
        sum += matrix[i * cols + j];
\end{lstlisting}

\section{False sharing i współbieżne liczniki}
\label{sec:cpp10-false-sharing}

\subsection{Na czym polega false sharing?}

Jeżeli dwa wątki wielokrotnie zapisują \emph{różne} zmienne, ale
zmienne znajdują się w tej samej linii cache, protokół spójności
pamięci podręcznej może wymuszać kosztowną wymianę własności tej
linii między rdzeniami. Nazywamy to \emph{false sharing}.
Nie oznacza to \emph{data race}, jeżeli poszczególne zmienne są
atomowe lub inaczej poprawnie synchronizowane.

\begin{lstlisting}
#include <atomic>
#include <cstdint>

struct NaiveCounters {
    std::atomic<std::uint64_t> first{0};
    std::atomic<std::uint64_t> second{0};
};
\end{lstlisting}

Dwa liczniki mogą być blisko siebie, jednak dokładny layout obiektu
i mapowanie do linii cache nie są gwarantowane. Na niektórych
platformach ich równoczesna modyfikacja będzie kosztowna.

\subsection{Rozdzielenie liczników w pamięci}

\begin{lstlisting}
#include <atomic>
#include <cstdint>

struct alignas(64) IsolatedAtomic {
    std::atomic<std::uint64_t> value{0};
};
struct SeparatedCounters {
    IsolatedAtomic first, second;
};
\end{lstlisting}

Tu użyliśmy 64 bajtów wyłącznie jako \emph{przykładowej} wartości.
Od C++17 w nagłówku \texttt{<new>} istnieją dwie stałe
opisujące zalecane odstępy dla danych często modyfikowanych przez
różne wątki (\emph{destructive interference}) oraz dla danych,
które powinny pozostawać blisko siebie (\emph{constructive
interference}). Pełną nazwę pierwszej z nich pokazuje kod poniżej.
Wartości stałych zależą od implementacji; nie stanowią pomiaru
rozmiaru linii cache w uruchomionym procesorze.

\begin{lstlisting}
#include <atomic>
#include <new>

struct alignas(
    std::hardware_destructive_interference_size
) Counter {
    std::atomic<unsigned> value{0};
};
\end{lstlisting}

Rozdzielenie liczników zwiększa rozmiar struktury. Przy milionie
liczników koszt pamięciowy może przewyższyć korzyść. Technika ma
największy sens, gdy oddzielne wątki często zapisują niewielką
liczbę niezależnych pól.

\subsection{False sharing a semantyka atomowa}

\texttt{std::memory\_order\_relaxed} wystarczy do niektórych
niezależnych liczników, gdy zależy nam wyłącznie na atomowym
zliczaniu; nie gwarantuje jednak synchronizacji innych danych.
Usunięcie false sharing nie zastępuje poprawnego modelu pamięci
współbieżnej. Szczegółowe zasady operacji atomowych były omawiane
w rozdziale~1 o współbieżności.

\section{Jak rzetelnie mierzyć wydajność?}
\label{sec:cpp10-benchmark}

\subsection{Zasady pomiarów}

Najważniejsze pytanie brzmi nie „czy ten kod wygląda szybciej?”,
lecz „czy w rzeczywistym scenariuszu uzyskujemy powtarzalną poprawę?”.
Doświadczenie powinno uwzględniać:
\begin{enumerate}
\item kompilację zoptymalizowaną (np. \texttt{-O2} lub \texttt{-O3});
\item \emph{warm-up} i wielokrotne uruchomienia;
\item jednakowe dane oraz taki sam rezultat dla obu wariantów;
\item ochronę przed usunięciem testowanej pracy przez optymalizator;
\item pomiar rozkładu wyników, a nie pojedynczego odczytu;
\item ograniczenie wpływu alokacji, I/O, schedulerów i skalowania CPU;
\item kontrolę zużycia pamięci, a nie tylko czasu;
\item narzędzia profilujące (np. Linux \texttt{perf}, odpowiedniki
      sprzętowe lub profilery danego środowiska).
\end{enumerate}

\subsection{Minimalna funkcja pomiarowa}

\begin{lstlisting}
#include <chrono>
#include <utility>

template<class F>
auto measure(F&& f) {
    auto start = std::chrono::steady_clock::now();
    auto result = std::forward<F>(f)();
    auto stop = std::chrono::steady_clock::now();
    return std::pair{result, stop - start};
}
\end{lstlisting}

Taka funkcja przydaje się do nauki, ale \emph{nie jest kompletna}
biblioteka benchmarkowa. Dla zbyt krótkich fragmentów narzut
pomiaru może przeważać nad czasem pracy. Dla liczb zmiennoprzecinkowych
porównywanie różnych kolejności sumowania wymaga ostrożności,
ponieważ zmiana kolejności działań może zmienić zaokrąglenia.

\subsection{Sanitizatory i wydajność}

Sanitizatory, np. AddressSanitizer i UndefinedBehaviorSanitizer,
są narzędziami do badania poprawności pamięci i innych błędów.
Wprowadzają narzut, dlatego \emph{pomiarów porównawczych wydajności}
nie należy bezrefleksyjnie wykonywać na buildach sanitizowanych.
Bezpieczny cykl pracy wygląda następująco: test poprawności,
profilowanie zwykłego buildu, zmiana jednej rzeczy, ponowny test
poprawności i ponowny pomiar.

\section{Przykładowy projekt: szybkie przetwarzanie wielu punktów}
\label{sec:cpp10-project}

Przyjmijmy, że program przechowuje milion punktów \((x,y,z,I)\)
i dla każdego oblicza tylko prosty filtr zależny od \(x\) i \(I\).
Możemy rozważyć dwa modele:
\begin{enumerate}
\item \textbf{AoS:} \texttt{vector<Point>} --- prostszy interfejs,
      wygodny dostęp do pełnego punktu.
\item \textbf{SoA:} osobne \texttt{vector<float>} na \(x,y,z,I\) ---
      możliwy mniejszy transfer danych potrzebnych do filtru.
\end{enumerate}

Następnie oceniamy, czy warto użyć \texttt{reserve} przed zapełnieniem
kontenerów; jeśli rekordy są tworzone na potrzeby jednego zapytania,
możemy sprawdzić arenę monotoniczną. Trzeba jednak rozróżnić
\emph{koszt wypełnienia struktur} od \emph{kosztu przetwarzania}: arena
zmienia zachowanie alokacji, nie gwarantuje poprawy pracy procesora
nad danymi.

Dla kodu numerycznego duży wpływ może mieć też wybór kolejności
wymiarów macierzy, reprezentacji sąsiedztwa w siatce czy sposobu
obliczania operatora na konkretnym punkcie zamiast konstruowania
dużych obiektów pośrednich. Właściwa reprezentacja danych potrafi
dać większy zysk niż drobne optymalizacje pojedynczych instrukcji.

\section{Jak dobrać technikę?}

\begin{center}
\small
\begin{tabular}{|p{5.7cm}|p{6.7cm}|}
\hline
\textbf{Problem} & \textbf{Pierwszy sensowny krok} \\
\hline
Za dużo dynamicznych alokacji & \texttt{reserve}, ograniczenie obiektów pośrednich, pomiar alokacji \\
\hline
Wiele obiektów znika razem & \texttt{pmr::monotonic\_buffer\_resource} po kontroli czasu życia \\
\hline
Dużo małych alokacji o różnej żywotności & \texttt{pmr::...pool\_resource} i pomiar \\
\hline
Praca nad jedną kolumną wielu rekordów & Porównanie AoS z SoA \\
\hline
Losowe skoki po pamięci & Sprawdzenie lokalności, zmian formatu danych i blokowania pętli \\
\hline
Niezależne atomiki spowalniają się nawzajem & Analiza false sharing i separacja intensywnie zapisywanych pól \\
\hline
Trzeba uzyskać silne wyrównanie & \texttt{alignas}, odpowiednia alokacja i kontrola \texttt{alignof} \\
\hline
Trzeba napisać kontener lub pulę & \texttt{allocator\_traits}, zarządzanie czasem życia, RAII \\
\hline
Potrzeba przenośnej serializacji & Jawny format bajtowy, bez kopiowania całych struktur \\
\hline
\end{tabular}
\end{center}

\section{Typowe błędy i pułapki}

\begin{enumerate}
\item \textbf{Wniosek „\texttt{sizeof(T)} to suma pól”.}
      Padding może zwiększyć rozmiar, a reguły klas mogą być złożone.
\item \textbf{Uznanie layoutu z jednej kompilacji za standard.}
      ABI, wyrównanie i rozmiary typów różnią się między platformami.
\item \textbf{Nieprawidłowe \texttt{reinterpret\_cast}.}
      Sama zmiana typu wskaźnika nie legalizuje dostępu do danych.
\item \textbf{Mylenie pamięci z czasem życia obiektu.}
      \texttt{allocate} nie zastępuje konstruowania obiektów klasowych.
\item \textbf{Zwalnianie pamięci bez destruktora.}
      Dla nietrywialnych typów może wyciekać zasób zarządzany przez obiekt.
\item \textbf{Niezgodna para alokacja--dealokacja.}
      Szczególnie ważne przy over-alignment i ręcznej alokacji.
\item \textbf{Pozostawienie kontenera po zniszczeniu areny \texttt{pmr}.}
      Jego pamięć staje się niepoprawna; wynik to niezdefiniowane zachowanie.
\item \textbf{Wywołanie \texttt{release()} przy żywych obiektach.}
      Arena może unieważnić ich pamięć bez wywoływania destruktorów.
\item \textbf{Założenie, że \texttt{monotonic\_buffer\_resource} oszczędza
      pamięć.} Zwolnienie pojedynczych bloków zwykle nie następuje.
\item \textbf{Współdzielenie \texttt{unsynchronized\_pool\_resource}
      między wątkami bez ochrony.} Zasób nie jest thread-safe.
\item \textbf{Optymalizacja poprzez mniejszą strukturę bez pomiaru.}
      Mniejszy \texttt{sizeof} nie gwarantuje krótszego czasu wykonania.
\item \textbf{Rzekoma optymalizacja przez \texttt{volatile}.}
      \texttt{volatile} nie zapewnia synchronizacji między wątkami.
\item \textbf{Pomiar jednego przebiegu.} Scheduler, częstotliwości
      zegara i cache mogą silnie zniekształcić wynik.
\item \textbf{Użycie sanitizatora jako miary szybkości.}
      Narzut instrumentacji zmienia charakterystykę programu.
\end{enumerate}

\paragraph{Pytania kontrolne.}
\begin{enumerate}
\item Czym różnią się \texttt{sizeof(T)} i \texttt{alignof(T)}?
\item Dlaczego zmiana kolejności pól może zmienić rozmiar struktury?
\item Czy placement \texttt{new} przydziela pamięć?
\item Jaka jest różnica między \texttt{destroy\_at} a \texttt{deallocate}?
\item Kiedy arena monotoniczna może zużywać więcej pamięci niż oczekujemy?
\item Dlaczego \texttt{pmr::vector<T>} nie jest właścicielem zasobu?
\item Kiedy SoA może przyspieszyć obliczenia na rekordach?
\item Czy false sharing i data race są tym samym problemem?
\item Dlaczego \texttt{std::memory\_order\_relaxed} nie rozwiązuje
      wszystkich problemów synchronizacji?
\item Jak odróżnić realną poprawę wydajności od szumu pomiarowego?
\end{enumerate}

\paragraph{Dalsza lektura.}
Dokumentacja \emph{cppreference} obejmuje hasła \emph{Object},
\emph{Object lifetime}, \texttt{alignof}, \texttt{alignas},
\texttt{std::construct\_at} oraz \texttt{allocator\_traits}.
Osobno warto przeczytać opis zasobów \texttt{std::pmr}
i własności \emph{hardware interference size}.
W praktyce polecam analizę śladów alokacji, profilery
wydajności i \emph{Compiler Explorer} do porównywania
wygenerowanego kodu maszynowego.
